\documentclass[12pt,a4paper]{article}
\usepackage[margin=2.4cm]{geometry}
\usepackage{amsmath,amssymb,bm,mathtools}
\usepackage{fontspec}
\setmainfont{TeX Gyre Termes}
\setsansfont{TeX Gyre Heros}
\setmonofont{Latin Modern Mono}[Scale=0.9]
\usepackage{graphicx,float,caption,subcaption,booktabs,longtable,array,multirow,tabularx}
\usepackage{xcolor}
\usepackage{setspace}
\usepackage{lineno}
\usepackage{microtype}
\usepackage{xurl}
\usepackage{siunitx}
\usepackage{enumitem}
\usepackage{fancyhdr}
\usepackage{titlesec}
\titleformat{\paragraph}[runin]{\normalfont\itshape}{}{0pt}{}
\titlespacing*{\paragraph}{0pt}{1.1ex plus .2ex}{0.7em}
\usepackage[authoryear,round,semicolon,sort&compress]{natbib}
\definecolor{lk}{RGB}{0,70,140}
\usepackage[colorlinks=true,linkcolor=lk,citecolor=lk,urlcolor=lk,breaklinks=true,pdfborder={0 0 0}]{hyperref}
\usepackage[nameinlink,capitalise]{cleveref}
\sisetup{detect-all,separate-uncertainty=false}
\setstretch{1.5}
\setlength{\parskip}{3pt}
\captionsetup{font=small,labelfont=bf,justification=centering,singlelinecheck=true}
\titleformat{\section}{\large\bfseries}{\thesection.}{0.5em}{}
\titleformat{\subsection}{\normalsize\bfseries}{\thesubsection}{0.5em}{}
\titleformat{\subsubsection}{\normalsize\itshape}{\thesubsubsection}{0.5em}{}
\pagestyle{fancy}\fancyhf{}\fancyfoot[C]{\small\thepage}\renewcommand{\headrulewidth}{0pt}
\setlength{\headheight}{14pt}
\bibliographystyle{plainnathl}
\newcommand{\doihl}[1]{doi: \href{https://doi.org/#1}{\nolinkurl{#1}}}
\newcommand{\dd}{\mathrm{d}}
\newcommand{\ET}{\mathit{ET}}
\newcommand{\Ent}{\mathit{ET}_{\mathrm{nt}}}
\newcommand{\Ks}{K_{\mathrm{s}}}
\newcommand{\Se}{S_{\mathrm{e}}}
\newcommand{\TAW}{\mathit{TAW}}
\newcommand{\repourl}{https://github.com/halilunaziru73-creator/Coupled-Water-Fluxes-Pore-Scale-Dynamics-Precision-Irrigation-Mulching-Feedback-Loops}
\newcommand{\repo}{\href{\repourl}{repository}}
\newcommand{\dataurl}{\repourl/tree/main/Data}
\newcommand{\codeurl}{\repourl/tree/main/Manuscripts}
\newcommand{\softurl}{\repourl/tree/main/Software}
\newcommand{\datalink}{\href{\dataurl}{Data}}
\newcommand{\codelink}{\href{\codeurl}{Code}}
\newcommand{\softlink}{\href{\softurl}{Software}}
\newcommand{\manuallink}{\href{\manualurl}{Manual}}
\newcommand{\downloadlink}{\href{\exeurl}{SWAFL.exe}}
\newcommand{\figdir}{../figures}
\newcommand{\manualurl}{https://github.com/halilunaziru73-creator/Coupled-Water-Fluxes-Pore-Scale-Dynamics-Precision-Irrigation-Mulching-Feedback-Loops/raw/main/Software/v1.1.0/SWAFL_Desktop/docs/manual_v1.1.0/SWAFL_Manual_v1.1.0.pdf}
\newcommand{\exeurl}{https://github.com/halilunaziru73-creator/Coupled-Water-Fluxes-Pore-Scale-Dynamics-Precision-Irrigation-Mulching-Feedback-Loops/releases/latest/download/SWAFL.exe}
\newcommand{\authorblock}{%
\begin{center}
{\large\textbf{Naziru Halilu}}$^{1,2,*}$\\[3pt]
{\small $^{1}$ Agricultural University of Athens, Athens, Greece\\
$^{2}$ Public University of Navarre, Pamplona, Spain\\
ORCID: \href{https://orcid.org/0009-0006-2324-9599}{0009-0006-2324-9599}\\
$^{*}$ Correspondence: \href{mailto:halilunaziru73@gmail.com}{halilunaziru73@gmail.com}}
\end{center}}
\newcommand{\csurl}{https://github.com/halilunaziru73-creator/Coupled-Water-Fluxes-Pore-Scale-Dynamics-Precision-Irrigation-Mulching-Feedback-Loops/tree/main/Data/Comprehensive_Summary}
\newcommand{\cslink}{\href{\csurl}{\nolinkurl{Data/Comprehensive_Summary}}}
\newcommand{\csavail}{\par\smallskip\textbf{Full simulation outputs.} Every value of every log of every run (the annual water and nitrogen balances, the harvests and the crop state, year by year), one compressed table per experiment (\texttt{E00\_Baselines} to \texttt{E16\_Edge\_Cases}), the one-row-per-run summary of all 10\,190 runs (\texttt{00\_summary\_all\_simulations}), the per-experiment summary (\texttt{00\_SUMMARY\_ALL\_EXPERIMENTS}) and the script that made them (\texttt{summarize\_results3.py}) are in the folder \cslink{} of the repository (\url{\csurl}); their 20-year means reproduce the run table within 2\,\%.}
% Data and code availability of the two analysis papers (P1, P2)
\newcommand{\dataavail}{%
\textbf{Data.} The complete Daisy ensemble outputs (all 10\,190 runs, raw and derived) are archived on Zenodo, \doihl{10.5281/zenodo.23043626}. The raw run table used here, and a linked list of every data file used, are in the \datalink{} folder of the project repository (\url{\dataurl}).\csavail
}
% Software, data and code availability of the software paper (P3)
\newcommand{\softavail}{%
\textbf{Software.} SWAFL Desktop (version 1.1.0) is open source; the \softlink{} source is in the project repository (\url{\softurl}) and the program for Windows is \downloadlink{} (\url{\exeurl}).
\par\smallskip\textbf{Manual.} The \manuallink{} (\url{\manualurl}) documents every control, view and equation and labels the provenance of all 134 numbered equations; it can be reviewed or downloaded as a PDF.
\par\smallskip\textbf{Data.} The complete Daisy ensemble outputs (all 10\,190 runs, raw and derived) are archived on Zenodo, \doihl{10.5281/zenodo.23043626}. The raw run table, and a linked list of every data file used, are in the \datalink{} folder of the repository (\url{\dataurl}).\csavail
}

\usepackage{pdflscape}
\setlength{\emergencystretch}{3em}

% A shared availability statement
\newcommand{\suiteavail}{%
\textbf{Software.} SWAFL Desktop (version 1.1.0) is open source; the \softlink{} source is in the project repository (\url{\softurl}) and the program for Windows is \downloadlink{} (\url{\exeurl}). The whole workbench also runs in a browser, with nothing to install (link in the repository README).
\par\smallskip\textbf{Manual.} The \manuallink{} (\url{\manualurl}) documents every control, screen and equation of the release, with labelled screenshots.
\par\smallskip\textbf{Data.} The complete Daisy ensemble outputs (all 10\,190 runs, raw and derived) are archived on Zenodo, \doihl{10.5281/zenodo.23043626}; the run table and a summary of every experiment are in the \datalink{} folder (\url{\dataurl}).\csavail
}

\renewcommand{\thesection}{S\arabic{section}}
\renewcommand{\theequation}{S\arabic{equation}}
\renewcommand{\thetable}{S\arabic{table}}
\renewcommand{\thefigure}{S\arabic{figure}}
\hypersetup{pdftitle={Supplementary Material: Soil Hydraulic Memory and the Recoverability of Evaporation Savings and the Closure of the Water Balance in a Simulation Study with Daisy},pdfauthor={Naziru Halilu}}
\begin{document}
\begin{center}{\large\bfseries Supplementary Material\par}\vspace{4pt}
{\Large\bfseries Soil Hydraulic Memory and the Recoverability of Evaporation Savings and the Closure of the Water Balance in a Simulation Study with Daisy\par}\vspace{6pt}
{\normalsize Derivations of the equations, the mechanisms behind the results, the data, and the origin of every table and figure\par}\end{center}
\authorblock
\tableofcontents
\clearpage

\section{Derivations of the equations}
\subsection{The retention and conductivity functions}
With the capillary law $r\propto1/h$, \citet{Mualem1976} expressed the relative conductivity of a bundle of connected capillaries as
\begin{equation}K_r(S_e)=S_e^{\lambda}\left[\int_0^{S_e}h(x)^{-1}\dd x\Big/\int_0^{1}h(x)^{-1}\dd x\right]^{2}.\end{equation}
Inserting the inverse of the van~Genuchten curve, $h(x)=\alpha^{-1}(x^{-1/m}-1)^{1/n}$, and substituting $y=x^{1/m}$ gives $\int_0^{S_e}h^{-1}\dd x=\alpha m\int_0^{S_e^{1/m}}y^{m-1+1/n}(1-y)^{-1/n}\dd y$. With $m=1-1/n$ the power of $y$ vanishes and the integral is $\alpha[1-(1-S_e^{1/m})^m]$; dividing by its value at $S_e=1$ ($\alpha$) and squaring yields the Mualem conductivity of the paper, $K=K_sS_e^{\lambda}[1-(1-S_e^{1/m})^m]^2$. For large $\alpha h$, $S_e\simeq(\alpha h)^{-(n-1)}$ and $K_r\simeq m^2(\alpha h)^{-[2n+\lambda(n-1)]}$: the conductivity of a soil with large $n$ falls faster with suction, which is the property that orders hydraulic memory.

\subsection{The water balance and the storage release}
Over a period, mass conservation of the soil column with surface runoff $Q_r$ is $P_r+I=\ET+D+\Delta S+Q_r$. Rearranging,
\begin{equation}\bar R\equiv\ET+D-P_r-I=-(\Delta S+Q_r),\end{equation}
so the storage release is positive when the column loses stored water (or sheds no runoff) and supplies the fluxes. In the initial-moisture experiment, a wetter initial state raises the stored water that is released over the run, which is the variation used to measure the memory fates.

\subsection{The memory fate as a pooled within-group slope}
Within a group of runs with the same rainfall scale $P$, write each flux as $x_c=a_P+\phi_x\bar R_c+e_c$. Minimising $\sum_P\sum_{c\in P}e_c^2$ over the group intercepts $a_P$ gives $a_P=\bar x_P-\phi_x\overline{\bar R}_P$, and then over $\phi_x$
\begin{equation}\phi_x=\frac{\sum_P\sum_{c\in P}(\bar R_c-\overline{\bar R}_P)(x_c-\bar x_P)}{\sum_P\sum_{c\in P}(\bar R_c-\overline{\bar R}_P)^2}.\end{equation}
Because the four fluxes add to $\ET+D=\bar R+P_r+I$ and $P_r+I$ is constant within a group, the deviations of the four fluxes add to the deviation of $\bar R$, and the four slopes add to one: $\phi_{E_s}+\phi_T+\phi_D+\phi_{\mathrm{oth}}=1$. Each $\phi_x$ is therefore the share of a released millimetre that leaves by flux $x$.

\subsection{Loss fates and the recoverable share}
Between two levels of the evaporation factor the change of the non-transpirational evapotranspiration $\delta\Ent$ must be balanced by the other fluxes. With storage changes negligible over 20 years, $\delta\Ent+\delta T+\delta D-\delta I\simeq0$, hence
\begin{equation}\psi_T+\psi_D+\psi_I\simeq1,\qquad \psi_x=\mp\frac{\delta x}{\delta\Ent},\end{equation}
the signs chosen so that each $\psi$ is positive when the flux compensates the evaporation change. Of a millimetre of evaporation saved, $\psi_D$ leaves as drainage and is not recovered by the crop or the irrigation account; the recoverable share is $\rho_R=1-\psi_D$.

\subsection{The rate-competition law}
Let stored water $S$ above field capacity leave by two parallel first-order sinks, drainage with rate constant $k_D\propto K_s/L$ and withdrawal by evaporation and uptake with $k_W\propto q^\ast/L$, where $L$ is a common length. Then $\dd S/\dd t=-(k_D+k_W)S$ and the drainage leaving over the release is
\begin{equation}\int_0^\infty k_DS\,\dd t=\frac{k_D}{k_D+k_W}S_0\quad\Longrightarrow\quad\phi_D=\frac{K_s}{K_s+q^\ast}.\end{equation}
The law is bounded in $[0,1]$, tends to 0 for an impermeable soil and to 1 for a free-draining one, and is the only one-parameter form with these limits that is linear in $K_s$ at small $K_s$. By the Buckingham theorem, $\phi_D=F(K_s/q^\ast,n,\alpha h^\ast)$; the law keeps the first group only.

\subsection{The invariant search by singular-value decomposition}
Arrange the 20-year mean fluxes of all runs as the rows of a matrix $\mathbf F$ (columns $P_r,I,\ET,D$, centred). A linear invariant is a vector $\mathbf c$ with $\mathbf F\mathbf c\approx\mathbf 0$. The best such vector in the least-squares sense, with $\|\mathbf c\|=1$, minimises $\|\mathbf F\mathbf c\|^2=\mathbf c^\top\mathbf F^\top\mathbf F\mathbf c$, which is the eigenvector of $\mathbf F^\top\mathbf F$ with the smallest eigenvalue, that is the right singular vector of $\mathbf F$ with the smallest singular value. Rescaled so that the $\ET$ coefficient is $-1$, it is the discovered balance; the permutation null shuffles each column independently and so destroys any true invariant, which gives the distribution of the smallest singular value under no structure.

\subsection{The closure residual and the unattributed term}
With $W=P_r+I$, the residual $\varepsilon=(W-\ET-D)/W\simeq(\Delta S+Q_r)/W$ is the share of the input not leaving by evapotranspiration or drainage. The unattributed evapotranspiration is $R=\ET-T-E_s-E_c$; in the full outputs it equals the evaporation of ponded water, litter and snow (data folder, closure within 1.1\,mm).

\subsection{Log-ratio coordinates and the sparse law}
Fluxes that share a total are compositional; the centred log-ratio $\mathrm{clr}(\mathbf x)_i=\ln x_i-\tfrac15\sum_j\ln x_j$ removes the constraint. The sparse laws regress $\ln y$ on a small set of log or dimensionless candidates chosen by forward selection; out-of-sample skill is $R^2_{\mathrm{cv}}=1-\sum(\ln y_i-\widehat{\ln y}_i^{(-g)})^2/\sum(\ln y_i-\overline{\ln y})^2$, with each prediction made by a model fitted without the group $g$ (experiment or soil) of run $i$.

\subsection{Irrigation compensation}
With the same balance, a change of rainfall $\Delta P_r$ at fixed trigger is offset partly by irrigation: $\kappa=-\Delta I/\Delta P_r$. When drainage and evapotranspiration do not change, $\kappa=1$; when part of the extra rain drains, $\kappa<1$.

\section{Why the results happen}
\paragraph{Why hydraulic memory differs between soils.} Water released from storage leaves by the faster of two routes. In coarse soils (large $n$, high $K_s$) the conductivity at field capacity is high, so released water drains before roots and the surface can use it; in fine soils the conductivity is low and the same water stays near the roots and the surface, where it is transpired or evaporated.
\paragraph{Why the memory transfers to other experiments.} The same competition decides where any extra water goes, whether it comes from initial storage, from a lower evaporation factor or from a mulch; the share is a property of the soil's conductivity, not of the source of the water.
\paragraph{Why warming raises and wetter climates lower the recoverable share.} Warming raises the evaporative demand, so the withdrawal side of the competition is faster and less of a saved millimetre drains. A wetter climate keeps the profile closer to field capacity, where drainage is fast, so more of a saving drains.
\paragraph{Why nitrogen has no effect.} The nitrogen rate changes the crop's growth but not the hydraulic properties, and in these runs it changes the water fluxes too little to move the shares; it therefore serves as a negative control.
\paragraph{Why an evaporative term is unattributed.} The reported components cover transpiration, soil evaporation and canopy evaporation; evaporation of ponded water, litter and snow is part of evapotranspiration but not of those components, so it appears as a residual in the run table.
\paragraph{Why evapotranspiration follows a compact law and drainage does not.} Evapotranspiration is limited by energy and by water and is smooth in its inputs over 20 years. Drainage is a threshold process (water drains only above field capacity) whose annual total depends on the timing of rain and irrigation events, which the long-term means do not contain.
\paragraph{Why judging a saving against soil evaporation alone misleads.} A mulch changes soil evaporation and also the unattributed evaporation; if only the reported soil evaporation is used as the denominator, the drainage share of the change is computed against the wrong total and can change sign.

\section{The data}
\Cref{tab:data} lists the 17 experiments of the data folder and \cref{fig:data} shows their outputs. The tables that follow are the result files of the analyses of the paper, printed in full or up to 220 rows; each is also in the workbook \nolinkurl{Data/Manuscript_Data/SWAFL_Manuscripts_Data.xlsx}, and every run is in \cslink.
\begin{table}[H]\centering\scriptsize
\caption{The data of this study: the 17 Daisy experiments of \texttt{Data/Comprehensive\_Summary} (10190 runs, 10041 valid; 20 years, 1980--1999, of Taastrup weather). Means of the valid runs; $Y$ yield (Mg\,ha$^{-1}$), $I$ irrigation, $D$ drainage at 1\,m and $\ET$ (mm\,yr$^{-1}$), $N_L$ nitrate leached (kg\,N\,ha$^{-1}$\,yr$^{-1}$), $t$ run time of Daisy (s). Every row is one file of the folder, linked.}
\label{tab:data}
\setlength{\tabcolsep}{3pt}
\begin{tabular}{@{}llrrp{3.5cm}rrrrrr@{}}
\toprule
Exp. & File & Runs & Valid & Factors varied & $Y$ & $I$ & $D$ & $\ET$ & $N_L$ & $t$\\\midrule
E0 & \href{\repourl/blob/main/Data/Comprehensive_Summary/E00_Baselines.csv.gz}{\texttt{Baselines}} & 10 & 10 & soil & 2.19 & 0 & 55 & 418 & 158 & 231\\
E1 & \href{\repourl/blob/main/Data/Comprehensive_Summary/E01_Doses.csv.gz}{\texttt{Doses}} & 200 & 189 & soil, rate, dose & 5.13 & 198 & 98 & 564 & 79 & 195\\
E2 & \href{\repourl/blob/main/Data/Comprehensive_Summary/E02_TriggerPressure.csv.gz}{\texttt{TriggerPressure}} & 200 & 197 & soil, trigger & 4.70 & 169 & 80 & 559 & 90 & 139\\
E3 & \href{\repourl/blob/main/Data/Comprehensive_Summary/E03_SensorDepth.csv.gz}{\texttt{SensorDepth}} & 200 & 200 & soil, sensor depth & 4.94 & 199 & 86 & 583 & 86 & 146\\
E4 & \href{\repourl/blob/main/Data/Comprehensive_Summary/E04_Methods.csv.gz}{\texttt{Methods}} & 100 & 96 & soil, method & 5.01 & 164 & 88 & 542 & 83 & 147\\
E5 & \href{\repourl/blob/main/Data/Comprehensive_Summary/E05_SHP_Dynamics.csv.gz}{\texttt{SHP Dynamics}} & 300 & 300 & soil, method, tillage & 4.13 & 152 & 43 & 579 & 84 & 189\\
E6 & \href{\repourl/blob/main/Data/Comprehensive_Summary/E06_Mediterranean_Pollution.csv.gz}{\texttt{Mediterranean Pollution}} & 210 & 193 & soil, method, trigger, N rate & 3.64 & 84 & 239 & 620 & 62 & 118\\
E7 & \href{\repourl/blob/main/Data/Comprehensive_Summary/E07_Mulching.csv.gz}{\texttt{Mulching}} & 200 & 200 & soil, mulch & 5.29 & 183 & 96 & 555 & 74 & 161\\
E8 & \href{\repourl/blob/main/Data/Comprehensive_Summary/E08_Rain_Intensity_Moisture.csv.gz}{\texttt{Rain Intensity Moisture}} & 1000 & 1000 & initial suction, rain scale & 0.68 & 0 & -1 & 282 & -0 & 6\\
E9 & \href{\repourl/blob/main/Data/Comprehensive_Summary/E09_Soil_Moisture_Rain.csv.gz}{\texttt{Soil Moisture Rain}} & 1000 & 987 & soil, initial suction, rain scale & 0.75 & 0 & 19 & 247 & 33 & 7\\
E10 & \href{\repourl/blob/main/Data/Comprehensive_Summary/E10_Amount_Timing_Soil.csv.gz}{\texttt{Amount Timing Soil}} & 1000 & 954 & soil, rate, dose, window & 4.42 & 139 & 76 & 532 & 97 & 146\\
E11 & \href{\repourl/blob/main/Data/Comprehensive_Summary/E11_Temp_Evap_Rain.csv.gz}{\texttt{Temp Evap Rain}} & 1000 & 1000 & rain scale, warming, demand & 3.81 & 324 & 18 & 754 & 66 & 151\\
E12 & \href{\repourl/blob/main/Data/Comprehensive_Summary/E12_Sensor_Trigger_Rain.csv.gz}{\texttt{Sensor Trigger Rain}} & 1000 & 1000 & trigger, sensor depth, rain scale & 5.88 & 260 & 35 & 670 & 32 & 155\\
E13 & \href{\repourl/blob/main/Data/Comprehensive_Summary/E13_N_Water_Soil.csv.gz}{\texttt{N Water Soil}} & 1000 & 975 & soil, method, rate, dose, trigger, N rate & 4.00 & 162 & 83 & 547 & 79 & 207\\
E14 & \href{\repourl/blob/main/Data/Comprehensive_Summary/E14_Mulch_Rain.csv.gz}{\texttt{Mulch Rain}} & 1000 & 1000 & rain scale, mulch & 6.15 & 235 & 43 & 634 & 25 & 173\\
E15 & \href{\repourl/blob/main/Data/Comprehensive_Summary/E15_Method_Depth_Amount.csv.gz}{\texttt{Method Depth Amount}} & 1000 & 1000 & method, rate, dose, root depth & 6.02 & 217 & 61 & 619 & 37 & 157\\
E16 & \href{\repourl/blob/main/Data/Comprehensive_Summary/E16_Edge_Cases.csv.gz}{\texttt{Edge Cases}} & 770 & 740 & soil, method, rate, dose, trigger, sensor depth, rain scale, warming, N rate, demand & 2.17 & 344 & 185 & 600 & 93 & 224\\
\bottomrule
\end{tabular}
\end{table}

\begin{figure}[H]\centering
\includegraphics[width=\textwidth]{\figdir/data/fig_data_overview.pdf}
\caption{Yield, irrigation, drainage and nitrate leaching of the valid runs, by experiment.}\label{fig:data}
\end{figure}
\subsection*{p1\_audit\_link.json}
\noindent\small Source: \nolinkurl{Manuscripts/results/p1_audit_link.json}; 46 values. Workbook sheet \texttt{R\_p1\_audit\_link}.\par
{\scriptsize\begin{longtable}{p{\dimexpr 0.4557\linewidth-2\tabcolsep\relax}p{\dimexpr 0.5443\linewidth-2\tabcolsep\relax}}\toprule\textbf{quantity} & \textbf{value}\\\midrule\endhead
{}E7.S01.psiD\_Ent & 0.2975\\
{}E7.S01.psiD\_Es & 0.3089\\
{}E7.S01.share\_of\_dEnt\_in\_Es & 0.5618\\
{}E7.S01.n & 20\\
{}E7.S02.psiD\_Ent & 0.6073\\
{}E7.S02.psiD\_Es & 0.3798\\
{}E7.S02.share\_of\_dEnt\_in\_Es & 0.4582\\
{}E7.S02.n & 20\\
{}E7.S03.psiD\_Ent & 0.7062\\
{}E7.S03.psiD\_Es & -0.9307\\
{}E7.S03.share\_of\_dEnt\_in\_Es & -0.3733\\
{}E7.S03.n & 20\\
{}E7.S04.psiD\_Ent & 0.6867\\
{}E7.S04.psiD\_Es & -0.3497\\
{}E7.S04.share\_of\_dEnt\_in\_Es & -0.3967\\
{}E7.S04.n & 20\\
{}E7.S05.psiD\_Ent & 0.6628\\
{}E7.S05.psiD\_Es & 0.157\\
{}E7.S05.share\_of\_dEnt\_in\_Es & 0.06688\\
{}E7.S05.n & 20\\
{}E7.S06.psiD\_Ent & 0.5212\\
{}E7.S06.psiD\_Es & 0.4076\\
{}E7.S06.share\_of\_dEnt\_in\_Es & 0.4787\\
{}E7.S06.n & 20\\
{}E7.S07.psiD\_Ent & 0.3955\\
{}E7.S07.psiD\_Es & 0.4209\\
{}E7.S07.share\_of\_dEnt\_in\_Es & 0.5235\\
{}E7.S07.n & 20\\
{}E7.S08.psiD\_Ent & 0.5489\\
{}E7.S08.psiD\_Es & 0.359\\
{}E7.S08.share\_of\_dEnt\_in\_Es & 0.3922\\
{}E7.S08.n & 20\\
{}E7.S09.psiD\_Ent & 0.4257\\
{}E7.S09.psiD\_Es & 0.441\\
{}E7.S09.share\_of\_dEnt\_in\_Es & 0.5036\\
{}E7.S09.n & 20\\
{}E7.S10.psiD\_Ent & 0.4137\\
{}E7.S10.psiD\_Es & 0.4541\\
{}E7.S10.share\_of\_dEnt\_in\_Es & 0.4989\\
{}E7.S10.n & 20\\
{}E7\_summary.median\_share\_of\_dEnt\_in\_Es & 0.4684\\
{}E7\_summary.median\_abs\_diff\_psiD & 0.1517\\
{}E7\_summary.max\_abs\_diff\_psiD & 1.637\\
{}E7\_summary.n\_soils & 10\\
{}E7\_summary.n\_soils\_diff\_gt\_0p1 & 6\\
{}E7\_summary.n\_soils\_sign\_flip & 2\\
\bottomrule\end{longtable}}

\subsection*{p1\_core.json}
\noindent\small Source: \nolinkurl{Manuscripts/results/p1_core.json}; 270 values, first 220 shown. Workbook sheet \texttt{R\_p1\_core}.\par
{\scriptsize\begin{longtable}{p{\dimexpr 0.4557\linewidth-2\tabcolsep\relax}p{\dimexpr 0.5443\linewidth-2\tabcolsep\relax}}\toprule\textbf{quantity} & \textbf{value}\\\midrule\endhead
{}memory.S01.fE & 0.1451\\
{}memory.S01.fT & 0.75\\
{}memory.S01.fD & 0.02912\\
{}memory.S01.fO & 0.07576\\
{}memory.S01.ci.soil\_evap\_mm & 0.113, 0.2002\\
{}memory.S01.ci.transpiration\_mm & 0.7102, 0.7721\\
{}memory.S01.ci.drainage\_1m\_mm & 0.02449, 0.03256\\
{}memory.S01.ci.oth & 0.06013, 0.09077\\
{}memory.S01.capFC & 82.79\\
{}memory.S01.cap60 & 112.6\\
{}memory.S01.n & 100\\
{}memory.S02.fE & 0.1927\\
{}memory.S02.fT & 0.4187\\
{}memory.S02.fD & 0.265\\
{}memory.S02.fO & 0.1236\\
{}memory.S02.ci.soil\_evap\_mm & 0.1604, 0.2324\\
{}memory.S02.ci.transpiration\_mm & 0.4009, 0.4444\\
{}memory.S02.ci.drainage\_1m\_mm & 0.1982, 0.313\\
{}memory.S02.ci.oth & 0.1086, 0.1398\\
{}memory.S02.capFC & 104\\
{}memory.S02.cap60 & 145.7\\
{}memory.S02.n & 100\\
{}memory.S03.fE & 0.07403\\
{}memory.S03.fT & 0.07982\\
{}memory.S03.fD & 0.8455\\
{}memory.S03.fO & 0.0006454\\
{}memory.S03.ci.soil\_evap\_mm & 0.06681, 0.08221\\
{}memory.S03.ci.transpiration\_mm & 0.04333, 0.1216\\
{}memory.S03.ci.drainage\_1m\_mm & 0.7959, 0.89\\
{}memory.S03.ci.oth & 0.0002684, 0.001096\\
{}memory.S03.capFC & 4.172\\
{}memory.S03.cap60 & 9.87\\
{}memory.S03.n & 87\\
{}memory.S04.fE & 0.09992\\
{}memory.S04.fT & 0.1535\\
{}memory.S04.fD & 0.7409\\
{}memory.S04.fO & 0.005727\\
{}memory.S04.ci.soil\_evap\_mm & 0.093, 0.1086\\
{}memory.S04.ci.transpiration\_mm & 0.1329, 0.1871\\
{}memory.S04.ci.drainage\_1m\_mm & 0.7008, 0.7655\\
{}memory.S04.ci.oth & 0.004453, 0.007196\\
{}memory.S04.capFC & 13.19\\
{}memory.S04.cap60 & 25.76\\
{}memory.S04.n & 100\\
{}memory.S05.fE & 0.1277\\
{}memory.S05.fT & 0.2043\\
{}memory.S05.fD & 0.643\\
{}memory.S05.fO & 0.02492\\
{}memory.S05.ci.soil\_evap\_mm & 0.1174, 0.1419\\
{}memory.S05.ci.transpiration\_mm & 0.1839, 0.2358\\
{}memory.S05.ci.drainage\_1m\_mm & 0.5987, 0.6723\\
{}memory.S05.ci.oth & 0.02044, 0.0298\\
{}memory.S05.capFC & 41.33\\
{}memory.S05.cap60 & 71.37\\
{}memory.S05.n & 100\\
{}memory.S06.fE & 0.04019\\
{}memory.S06.fT & 0.5825\\
{}memory.S06.fD & 0.2892\\
{}memory.S06.fO & 0.08811\\
{}memory.S06.ci.soil\_evap\_mm & 0.02569, 0.05774\\
{}memory.S06.ci.transpiration\_mm & 0.554, 0.6231\\
{}memory.S06.ci.drainage\_1m\_mm & 0.239, 0.3242\\
{}memory.S06.ci.oth & 0.07656, 0.1007\\
{}memory.S06.capFC & 131.8\\
{}memory.S06.cap60 & 183.1\\
{}memory.S06.n & 100\\
{}memory.S07.fE & 0.04151\\
{}memory.S07.fT & 0.7155\\
{}memory.S07.fD & 0.1473\\
{}memory.S07.fO & 0.09566\\
{}memory.S07.ci.soil\_evap\_mm & 0.0304, 0.05195\\
{}memory.S07.ci.transpiration\_mm & 0.6886, 0.7481\\
{}memory.S07.ci.drainage\_1m\_mm & 0.1124, 0.1754\\
{}memory.S07.ci.oth & 0.0862, 0.1052\\
{}memory.S07.capFC & 208.1\\
{}memory.S07.cap60 & 248.9\\
{}memory.S07.n & 100\\
{}memory.S08.fE & 0.07846\\
{}memory.S08.fT & 0.5547\\
{}memory.S08.fD & 0.2852\\
{}memory.S08.fO & 0.08163\\
{}memory.S08.ci.soil\_evap\_mm & 0.06028, 0.1076\\
{}memory.S08.ci.transpiration\_mm & 0.5128, 0.593\\
{}memory.S08.ci.drainage\_1m\_mm & 0.2427, 0.3188\\
{}memory.S08.ci.oth & 0.07191, 0.09251\\
{}memory.S08.capFC & 57.04\\
{}memory.S08.cap60 & 84.54\\
{}memory.S08.n & 100\\
{}memory.S09.fE & 0.02717\\
{}memory.S09.fT & 0.7744\\
{}memory.S09.fD & 0.08702\\
{}memory.S09.fO & 0.1114\\
{}memory.S09.ci.soil\_evap\_mm & 0.01397, 0.03822\\
{}memory.S09.ci.transpiration\_mm & 0.7541, 0.7978\\
{}memory.S09.ci.drainage\_1m\_mm & 0.06404, 0.1082\\
{}memory.S09.ci.oth & 0.1008, 0.1222\\
{}memory.S09.capFC & 167.5\\
{}memory.S09.cap60 & 194.4\\
{}memory.S09.n & 100\\
{}memory.S10.fE & 0.02489\\
{}memory.S10.fT & 0.8416\\
{}memory.S10.fD & 0.02197\\
{}memory.S10.fO & 0.1115\\
{}memory.S10.ci.soil\_evap\_mm & 0.01344, 0.0364\\
{}memory.S10.ci.transpiration\_mm & 0.829, 0.8548\\
{}memory.S10.ci.drainage\_1m\_mm & 0.01336, 0.03072\\
{}memory.S10.ci.oth & 0.1015, 0.1211\\
{}memory.S10.capFC & 173\\
{}memory.S10.cap60 & 190.1\\
{}memory.S10.n & 100\\
{}rec16.S01.psiD & 0.2946\\
{}rec16.S01.psiI & 0.6189\\
{}rec16.S01.psiT & 0.1113\\
{}rec16.S01.psiS & 0.02473\\
{}rec16.S01.dE & 218.5\\
{}rec16.S01.ci.psiD & 0.1943, 0.3918\\
{}rec16.S01.ci.psiI & 0.448, 0.7595\\
{}rec16.S01.ci.psiT & -0.04115, 0.3197\\
{}rec16.S01.ci.psiS & 0.00506, 0.04393\\
{}rec16.S01.ci.dE & 179.2, 257.9\\
{}rec16.S01.n & 32\\
{}rec16.S02.psiD & 0.193\\
{}rec16.S02.psiI & 0.6578\\
{}rec16.S02.psiT & 0.157\\
{}rec16.S02.psiS & 0.007809\\
{}rec16.S02.dE & 143.6\\
{}rec16.S02.ci.psiD & 0.1389, 0.2439\\
{}rec16.S02.ci.psiI & 0.3244, 0.858\\
{}rec16.S02.ci.psiT & -0.02457, 0.4575\\
{}rec16.S02.ci.psiS & 0.004907, 0.0132\\
{}rec16.S02.ci.dE & 99.58, 201.3\\
{}rec16.S02.n & 13\\
{}rec16.S03.psiD & 1.21\\
{}rec16.S03.psiI & -0.3839\\
{}rec16.S03.psiT & 0.1785\\
{}rec16.S03.psiS & 0.004424\\
{}rec16.S03.dE & 99.64\\
{}rec16.S03.ci.psiD & 0.955, 1.46\\
{}rec16.S03.ci.psiI & -0.8056, -0.05112\\
{}rec16.S03.ci.psiT & 0.01523, 0.4111\\
{}rec16.S03.ci.psiS & -0.001305, 0.01098\\
{}rec16.S03.ci.dE & 86.72, 112.8\\
{}rec16.S03.n & 32\\
{}rec16.S04.psiD & 0.7993\\
{}rec16.S04.psiI & 0.2189\\
{}rec16.S04.psiT & -0.007351\\
{}rec16.S04.psiS & 0.01086\\
{}rec16.S04.dE & 115\\
{}rec16.S04.ci.psiD & 0.6691, 0.9158\\
{}rec16.S04.ci.psiI & 0.1352, 0.3027\\
{}rec16.S04.ci.psiT & -0.1054, 0.1089\\
{}rec16.S04.ci.psiS & 0.005401, 0.01805\\
{}rec16.S04.ci.dE & 99.12, 129.1\\
{}rec16.S04.n & 32\\
{}rec16.S05.psiD & 0.5821\\
{}rec16.S05.psiI & 0.4146\\
{}rec16.S05.psiT & 0.01543\\
{}rec16.S05.psiS & 0.0121\\
{}rec16.S05.dE & 141.1\\
{}rec16.S05.ci.psiD & 0.4566, 0.6966\\
{}rec16.S05.ci.psiI & 0.2391, 0.5755\\
{}rec16.S05.ci.psiT & -0.1053, 0.177\\
{}rec16.S05.ci.psiS & 0.003533, 0.02239\\
{}rec16.S05.ci.dE & 115.9, 167.7\\
{}rec16.S05.n & 32\\
{}rec16.S06.psiD & 0.469\\
{}rec16.S06.psiI & 0.4334\\
{}rec16.S06.psiT & 0.1068\\
{}rec16.S06.psiS & 0.009196\\
{}rec16.S06.dE & 191.5\\
{}rec16.S06.ci.psiD & 0.3626, 0.5992\\
{}rec16.S06.ci.psiI & 0.1834, 0.607\\
{}rec16.S06.ci.psiT & -0.0311, 0.3117\\
{}rec16.S06.ci.psiS & -0.003642, 0.02676\\
{}rec16.S06.ci.dE & 139.4, 247.3\\
{}rec16.S06.n & 32\\
{}rec16.S07.psiD & 0.3625\\
{}rec16.S07.psiI & 0.5278\\
{}rec16.S07.psiT & 0.1168\\
{}rec16.S07.psiS & 0.007042\\
{}rec16.S07.dE & 217.6\\
{}rec16.S07.ci.psiD & 0.2597, 0.4852\\
{}rec16.S07.ci.psiI & 0.2777, 0.7039\\
{}rec16.S07.ci.psiT & -0.03777, 0.348\\
{}rec16.S07.ci.psiS & -0.009443, 0.02719\\
{}rec16.S07.ci.dE & 168.3, 271.5\\
{}rec16.S07.n & 32\\
{}rec16.S08.psiD & 0.533\\
{}rec16.S08.psiI & 0.3857\\
{}rec16.S08.psiT & 0.09111\\
{}rec16.S08.psiS & 0.009891\\
{}rec16.S08.dE & 152.5\\
{}rec16.S08.ci.psiD & 0.4093, 0.6608\\
{}rec16.S08.ci.psiI & 0.1227, 0.5953\\
{}rec16.S08.ci.psiT & -0.07816, 0.3433\\
{}rec16.S08.ci.psiS & -0.005323, 0.02717\\
{}rec16.S08.ci.dE & 114.9, 190.4\\
{}rec16.S08.n & 32\\
{}rec16.S09.psiD & 0.3117\\
{}rec16.S09.psiI & 0.645\\
{}rec16.S09.psiT & 0.04777\\
{}rec16.S09.psiS & 0.004532\\
{}rec16.S09.dE & 217.4\\
{}rec16.S09.ci.psiD & 0.1979, 0.4445\\
{}rec16.S09.ci.psiI & 0.3863, 0.806\\
{}rec16.S09.ci.psiT & -0.07744, 0.2568\\
{}rec16.S09.ci.psiS & -0.01091, 0.02484\\
{}rec16.S09.ci.dE & 165.6, 267.7\\
{}rec16.S09.n & 32\\
{}rec16.S10.psiD & 0.3103\\
{}rec16.S10.psiI & 0.838\\
{}rec16.S10.psiT & -0.01707\\
{}rec16.S10.psiS & 0.1312\\
{}rec16.S10.dE & 222.1\\
{}rec16.S10.ci.psiD & 0.217, 0.4126\\
{}rec16.S10.ci.psiI & 0.6104, 1.033\\
{}rec16.S10.ci.psiT & -0.1507, 0.1387\\
{}rec16.S10.ci.psiS & 0.05407, 0.1948\\
{}rec16.S10.ci.dE & 176.8, 272.6\\
{}rec16.S10.n & 32\\
\bottomrule\end{longtable}}

\subsection*{p1\_ens\_structure.json}
\noindent\small Source: \nolinkurl{Manuscripts/results/p1_ens_structure.json}; 50 values. Workbook sheet \texttt{R\_p1\_ens\_structure}.\par
{\scriptsize\begin{longtable}{p{\dimexpr 0.4557\linewidth-2\tabcolsep\relax}p{\dimexpr 0.5443\linewidth-2\tabcolsep\relax}}\toprule\textbf{quantity} & \textbf{value}\\\midrule\endhead
{}dE.S01 & 54.3\\
{}dE.S02 & 28.7\\
{}dE.S03 & 3.4\\
{}dE.S04 & 6.9\\
{}dE.S05 & 17\\
{}dE.S06 & 32.8\\
{}dE.S07 & 45.3\\
{}dE.S08 & 25.5\\
{}dE.S09 & 41.2\\
{}dE.S10 & 44.8\\
{}dD.S01 & 16.5\\
{}dD.S02 & 20.8\\
{}dD.S03 & 6.1\\
{}dD.S04 & 7.9\\
{}dD.S05 & 9.5\\
{}dD.S06 & 19.5\\
{}dD.S07 & 25.4\\
{}dD.S08 & 15.4\\
{}dD.S09 & 25.8\\
{}dD.S10 & 27.1\\
{}dI.S01 & 46.6\\
{}dI.S02 & 10.4\\
{}dI.S03 & 0\\
{}dI.S04 & -0.4\\
{}dI.S05 & -5.5\\
{}dI.S06 & 18.1\\
{}dI.S07 & 45.1\\
{}dI.S08 & 9.2\\
{}dI.S09 & 42.8\\
{}dI.S10 & 49.1\\
{}tillage\_offset.S01 & 0.27\\
{}tillage\_offset.S02 & -1.82\\
{}tillage\_offset.S03 & 5.34\\
{}tillage\_offset.S04 & 3.69\\
{}tillage\_offset.S05 & 1.76\\
{}tillage\_offset.S06 & -1.32\\
{}tillage\_offset.S07 & -0.25\\
{}tillage\_offset.S08 & -3.69\\
{}tillage\_offset.S09 & -2.79\\
{}tillage\_offset.S10 & -1.7\\
{}tc\_effect.S01 & 0\\
{}tc\_effect.S02 & 0\\
{}tc\_effect.S03 & 0\\
{}tc\_effect.S04 & 0\\
{}tc\_effect.S05 & 0\\
{}tc\_effect.S06 & 1.09\\
{}tc\_effect.S07 & 0.23\\
{}tc\_effect.S08 & 2.54\\
{}tc\_effect.S09 & 2.06\\
{}tc\_effect.S10 & 1.48\\
\bottomrule\end{longtable}}

\subsection*{p1\_loso\_summary.json}
\noindent\small Source: \nolinkurl{Manuscripts/results/p1_loso_summary.json}; 5 values. Workbook sheet \texttt{R\_p1\_loso\_summary}.\par
{\scriptsize\begin{longtable}{p{\dimexpr 0.4557\linewidth-2\tabcolsep\relax}p{\dimexpr 0.5443\linewidth-2\tabcolsep\relax}}\toprule\textbf{quantity} & \textbf{value}\\\midrule\endhead
{}TAW & 0.1825\\
{}\$S\_e\$(100) & 0.1358\\
{}\$n\$ & 0.09716\\
{}\$/log K\_s\$* & 0.05811\\
{}mean only & 0.3178\\
\bottomrule\end{longtable}}

\subsection*{p1\_mech.json}
\noindent\small Source: \nolinkurl{Manuscripts/results/p1_mech.json}; 5 values. Workbook sheet \texttt{R\_p1\_mech}.\par
{\scriptsize\begin{longtable}{p{\dimexpr 0.4557\linewidth-2\tabcolsep\relax}p{\dimexpr 0.5443\linewidth-2\tabcolsep\relax}}\toprule\textbf{quantity} & \textbf{value}\\\midrule\endhead
{}hstar & 6.712\\
{}q & 15.41\\
{}pred & 0.8355, 0.7637, 0.6336, 0.3368, 0.1148, 0.2433, 0.08474, 0.02309\\
{}loo & 0.8219, 0.7791, 0.63, 0.3715, 0.1108, 0.2334, 0.08459, 0.02309\\
{}soils & S03, S04, S05, S06, S07, S08, S09, S10\\
\bottomrule\end{longtable}}

\subsection*{p1\_modifiers.json}
\noindent\small Source: \nolinkurl{Manuscripts/results/p1_modifiers.json}; 20 values. Workbook sheet \texttt{R\_p1\_modifiers}.\par
{\scriptsize\begin{longtable}{p{\dimexpr 0.4557\linewidth-2\tabcolsep\relax}p{\dimexpr 0.5443\linewidth-2\tabcolsep\relax}}\toprule\textbf{quantity} & \textbf{value}\\\midrule\endhead
{}fP.mean & 0.4193\\
{}fP.pos & 9\\
{}fP.n & 9\\
{}fP.p & 0.003906\\
{}fT.mean & -0.1612\\
{}fT.pos & 0\\
{}fT.n & 10\\
{}fT.p & 0.001953\\
{}fh.mean & 0.1364\\
{}fh.pos & 7\\
{}fh.n & 10\\
{}fh.p & 0.1602\\
{}fD.mean & -0.04023\\
{}fD.pos & 2\\
{}fD.n & 10\\
{}fD.p & 0.3223\\
{}fN.mean & -0.008285\\
{}fN.pos & 2\\
{}fN.n & 10\\
{}fN.p & 0.06445\\
\bottomrule\end{longtable}}

\subsection*{p1\_rank.json}
\noindent\small Source: \nolinkurl{Manuscripts/results/p1_rank.json}; 20 values. Workbook sheet \texttt{R\_p1\_rank}.\par
{\scriptsize\begin{longtable}{p{\dimexpr 0.4557\linewidth-2\tabcolsep\relax}p{\dimexpr 0.5443\linewidth-2\tabcolsep\relax}}\toprule\textbf{quantity} & \textbf{value}\\\midrule\endhead
{}\$n\$ vs \$/phi\_D\$.rho & 0.9515\\
{}\$n\$ vs \$/phi\_D\$.lo & 0.6883\\
{}\$n\$ vs \$/phi\_D\$.hi & 1\\
{}\$n\$ vs \$/phi\_D\$.p & 0.00015\\
{}\$/phi\_D\$ vs \$/psi\_D\$(E16).rho & 0.8545\\
{}\$/phi\_D\$ vs \$/psi\_D\$(E16).lo & 0.3703\\
{}\$/phi\_D\$ vs \$/psi\_D\$(E16).hi & 1\\
{}\$/phi\_D\$ vs \$/psi\_D\$(E16).p & 0.00195\\
{}\$/phi\_D\$ vs \$/psi\_D\$(E7).rho & 0.8909\\
{}\$/phi\_D\$ vs \$/psi\_D\$(E7).lo & 0.5472\\
{}\$/phi\_D\$ vs \$/psi\_D\$(E7).hi & 1\\
{}\$/phi\_D\$ vs \$/psi\_D\$(E7).p & 0.00135\\
{}\$n\$ vs \$/psi\_D\$(E16).rho & 0.697\\
{}\$n\$ vs \$/psi\_D\$(E16).lo & -0.0566\\
{}\$n\$ vs \$/psi\_D\$(E16).hi & 1\\
{}\$n\$ vs \$/psi\_D\$(E16).p & 0.03185\\
{}\$n\$ vs \$/psi\_D\$(E7).rho & 0.8909\\
{}\$n\$ vs \$/psi\_D\$(E7).lo & 0.4717\\
{}\$n\$ vs \$/psi\_D\$(E7).hi & 1\\
{}\$n\$ vs \$/psi\_D\$(E7).p & 0.001\\
\bottomrule\end{longtable}}

\subsection*{p1\_rate\_law.json}
\noindent\small Source: \nolinkurl{Manuscripts/results/p1_rate_law.json}; 17 values. Workbook sheet \texttt{R\_p1\_rate\_law}.\par
{\scriptsize\begin{longtable}{p{\dimexpr 0.4557\linewidth-2\tabcolsep\relax}p{\dimexpr 0.5443\linewidth-2\tabcolsep\relax}}\toprule\textbf{quantity} & \textbf{value}\\\midrule\endhead
{}q\_star\_cm\_d & 71.46\\
{}R2 & 0.9715\\
{}boot\_ci95 & 58.38, 105.2\\
{}n\_boot & 4999\\
{}jackknife\_range & 65.83, 79.97\\
{}loso\_rmse\_rate\_law & 0.05811\\
{}loso\_8.mean\_only & 0.3384\\
{}loso\_8.TAW & 0.08835\\
{}loso\_8.Se100 & 0.06407\\
{}loso\_8.n & 0.1004\\
{}loso\_8.log10Ks\_linear & 0.06183\\
{}loso\_10.mean\_only & 0.3178\\
{}loso\_10.TAW & 0.1825\\
{}loso\_10.Se100 & 0.1358\\
{}loso\_10.n & 0.09716\\
{}loo\_pred & 0.9116, 0.8418, 0.5702, 0.2502, 0.07804, 0.3136, 0.08007, 0.02297\\
{}soils & S03, S04, S05, S06, S07, S08, S09, S10\\
\bottomrule\end{longtable}}

\subsection*{p1\_table.csv}
\noindent\small Source: \nolinkurl{Manuscripts/results/p1_table.csv}; 10 rows. Workbook sheet \texttt{R\_p1\_table}.\par
{\scriptsize\setlength{\tabcolsep}{3pt}\begin{longtable}{>{\raggedright\arraybackslash}p{\dimexpr 0.1111\linewidth-2\tabcolsep\relax}>{\raggedright\arraybackslash}p{\dimexpr 0.1111\linewidth-2\tabcolsep\relax}>{\raggedright\arraybackslash}p{\dimexpr 0.1111\linewidth-2\tabcolsep\relax}>{\raggedright\arraybackslash}p{\dimexpr 0.1111\linewidth-2\tabcolsep\relax}>{\raggedright\arraybackslash}p{\dimexpr 0.1111\linewidth-2\tabcolsep\relax}>{\raggedright\arraybackslash}p{\dimexpr 0.1111\linewidth-2\tabcolsep\relax}>{\raggedright\arraybackslash}p{\dimexpr 0.1111\linewidth-2\tabcolsep\relax}>{\raggedright\arraybackslash}p{\dimexpr 0.1111\linewidth-2\tabcolsep\relax}>{\raggedright\arraybackslash}p{\dimexpr 0.1111\linewidth-2\tabcolsep\relax}}
\toprule
\textbf{Unnamed: 0} & \textbf{phiD\_\allowbreak{}E9} & \textbf{phiT\_\allowbreak{}E9} & \textbf{phiE\_\allowbreak{}E9} & \textbf{capFC} & \textbf{psiD\_\allowbreak{}E16} & \textbf{psiI\_\allowbreak{}E16} & \textbf{psiT\_\allowbreak{}E16} & \textbf{dE16}\\\midrule\endhead
{}S01 & 0.02912 & 0.75 & 0.1451 & 82.79 & 0.2946 & 0.6189 & 0.1113 & 218.5\\
{}S02 & 0.265 & 0.4187 & 0.1927 & 104 & 0.193 & 0.6578 & 0.157 & 143.6\\
{}S03 & 0.8455 & 0.07982 & 0.07403 & 4.172 & 1.21 & -0.3839 & 0.1785 & 99.64\\
{}S04 & 0.7409 & 0.1535 & 0.09992 & 13.19 & 0.7993 & 0.2189 & -0.007351 & 115\\
{}S05 & 0.643 & 0.2043 & 0.1277 & 41.33 & 0.5821 & 0.4146 & 0.01543 & 141.1\\
{}S06 & 0.2892 & 0.5825 & 0.04019 & 131.8 & 0.469 & 0.4334 & 0.1068 & 191.5\\
{}S07 & 0.1473 & 0.7155 & 0.04151 & 208.1 & 0.3625 & 0.5278 & 0.1168 & 217.6\\
{}S08 & 0.2852 & 0.5547 & 0.07846 & 57.04 & 0.533 & 0.3857 & 0.09111 & 152.5\\
{}S09 & 0.08702 & 0.7744 & 0.02717 & 167.5 & 0.3117 & 0.645 & 0.04777 & 217.4\\
{}S10 & 0.02197 & 0.8416 & 0.02489 & 173 & 0.3103 & 0.838 & -0.01707 & 222.1\\
\bottomrule
\end{longtable}}

\subsection*{p1\_table2.csv}
\noindent\small Source: \nolinkurl{Manuscripts/results/p1_table2.csv}; 10 rows. Workbook sheet \texttt{R\_p1\_table2}.\par
{\scriptsize\setlength{\tabcolsep}{3pt}\begin{longtable}{>{\raggedright\arraybackslash}p{\dimexpr 0.1111\linewidth-2\tabcolsep\relax}>{\raggedright\arraybackslash}p{\dimexpr 0.1111\linewidth-2\tabcolsep\relax}>{\raggedright\arraybackslash}p{\dimexpr 0.1111\linewidth-2\tabcolsep\relax}>{\raggedright\arraybackslash}p{\dimexpr 0.1111\linewidth-2\tabcolsep\relax}>{\raggedright\arraybackslash}p{\dimexpr 0.1111\linewidth-2\tabcolsep\relax}>{\raggedright\arraybackslash}p{\dimexpr 0.1111\linewidth-2\tabcolsep\relax}>{\raggedright\arraybackslash}p{\dimexpr 0.1111\linewidth-2\tabcolsep\relax}>{\raggedright\arraybackslash}p{\dimexpr 0.1111\linewidth-2\tabcolsep\relax}>{\raggedright\arraybackslash}p{\dimexpr 0.1111\linewidth-2\tabcolsep\relax}}
\toprule
\textbf{Unnamed: 0} & \textbf{phiD\_\allowbreak{}E9} & \textbf{phiT\_\allowbreak{}E9} & \textbf{phiE\_\allowbreak{}E9} & \textbf{capFC} & \textbf{psiD\_\allowbreak{}E16} & \textbf{psiI\_\allowbreak{}E16} & \textbf{psiT\_\allowbreak{}E16} & \textbf{dE16}\\\midrule\endhead
{}S01 & 0.02912 & 0.75 & 0.1451 & 82.79 & 0.2946 & 0.6189 & 0.1113 & 218.5\\
{}S02 & 0.265 & 0.4187 & 0.1927 & 104 & 0.193 & 0.6578 & 0.157 & 143.6\\
{}S03 & 0.8455 & 0.07982 & 0.07403 & 4.172 & 1.21 & -0.3839 & 0.1785 & 99.64\\
{}S04 & 0.7409 & 0.1535 & 0.09992 & 13.19 & 0.7993 & 0.2189 & -0.007351 & 115\\
{}S05 & 0.643 & 0.2043 & 0.1277 & 41.33 & 0.5821 & 0.4146 & 0.01543 & 141.1\\
{}S06 & 0.2892 & 0.5825 & 0.04019 & 131.8 & 0.469 & 0.4334 & 0.1068 & 191.5\\
{}S07 & 0.1473 & 0.7155 & 0.04151 & 208.1 & 0.3625 & 0.5278 & 0.1168 & 217.6\\
{}S08 & 0.2852 & 0.5547 & 0.07846 & 57.04 & 0.533 & 0.3857 & 0.09111 & 152.5\\
{}S09 & 0.08702 & 0.7744 & 0.02717 & 167.5 & 0.3117 & 0.645 & 0.04777 & 217.4\\
{}S10 & 0.02197 & 0.8416 & 0.02489 & 173 & 0.3103 & 0.838 & -0.01707 & 222.1\\
\bottomrule
\end{longtable}}

\subsection*{p2\_R.json}
\noindent\small Source: \nolinkurl{Manuscripts/results/p2_R.json}; 27 values. Workbook sheet \texttt{R\_p2\_R}.\par
{\scriptsize\begin{longtable}{p{\dimexpr 0.4557\linewidth-2\tabcolsep\relax}p{\dimexpr 0.5443\linewidth-2\tabcolsep\relax}}\toprule\textbf{quantity} & \textbf{value}\\\midrule\endhead
{}R\_model.terms & lnI, n, cap, dT\\
{}R\_model.coef.const & 3.355\\
{}R\_model.coef.lnI & 0.1944\\
{}R\_model.coef.n & 0.4917\\
{}R\_model.coef.cap & 0.1454\\
{}R\_model.coef.dT & 0.2383\\
{}R\_model.se.const & 0.02376\\
{}R\_model.se.lnI & 0.001882\\
{}R\_model.se.n & 0.01289\\
{}R\_model.se.cap & 0.004217\\
{}R\_model.se.dT & 0.00723\\
{}R\_model.r2 & 0.6111\\
{}R\_model.path[0] & lnI, 0.5509\\
{}R\_model.path[1] & n, 0.5356\\
{}R\_model.path[2] & cap, 0.5169\\
{}R\_model.path[3] & dT, 0.5137\\
{}R\_model.loeo\_rmse\_ln & 0.5137\\
{}R\_sparse.R2 & 0.3028\\
{}R\_sparse.rmse & 80.54\\
{}R\_sparse.R2ln & 0.4092\\
{}R\_GBM\_all.R2 & 0.4956\\
{}R\_GBM\_all.rmse & 68.51\\
{}R\_GBM\_all.R2ln & 0.4268\\
{}R\_GBM\_sel.R2 & 0.5095\\
{}R\_GBM\_sel.rmse & 67.56\\
{}R\_GBM\_sel.R2ln & 0.4665\\
{}R\_sparse\_LOSO\_R2ln & 0.4275\\
\bottomrule\end{longtable}}

\subsection*{p2\_a.json}
\noindent\small Source: \nolinkurl{Manuscripts/results/p2_a.json}; 46 values. Workbook sheet \texttt{R\_p2\_a}.\par
{\scriptsize\begin{longtable}{p{\dimexpr 0.4557\linewidth-2\tabcolsep\relax}p{\dimexpr 0.5443\linewidth-2\tabcolsep\relax}}\toprule\textbf{quantity} & \textbf{value}\\\midrule\endhead
{}svd.sv & 240.3, 82.6, 57.95, 39.37, 21.84, 18.17, 5.461\\
{}svd.null[0].k & 1\\
{}svd.null[0].sv & 5.461\\
{}svd.null[0].vec.rain\_mm & 0.8956\\
{}svd.null[0].vec.irrigation\_mm & 0.8343\\
{}svd.null[0].vec.ET\_mm & -1\\
{}svd.null[0].vec.transpiration\_mm & 0.3675\\
{}svd.null[0].vec.soil\_evap\_mm & 0.1207\\
{}svd.null[0].vec.canopy\_evap\_mm & -0.1216\\
{}svd.null[0].vec.drainage\_1m\_mm & -0.9325\\
{}svd.null[1].k & 2\\
{}svd.null[1].sv & 18.17\\
{}svd.null[1].vec.rain\_mm & -0.3585\\
{}svd.null[1].vec.irrigation\_mm & -0.2111\\
{}svd.null[1].vec.ET\_mm & -0.1925\\
{}svd.null[1].vec.transpiration\_mm & 0.7876\\
{}svd.null[1].vec.soil\_evap\_mm & 0.4982\\
{}svd.null[1].vec.canopy\_evap\_mm & 1\\
{}svd.null[1].vec.drainage\_1m\_mm & 0.1799\\
{}svd.null[2].k & 3\\
{}svd.null[2].sv & 21.84\\
{}svd.null[2].vec.rain\_mm & -0.03927\\
{}svd.null[2].vec.irrigation\_mm & 0.03899\\
{}svd.null[2].vec.ET\_mm & -0.02946\\
{}svd.null[2].vec.transpiration\_mm & -0.3726\\
{}svd.null[2].vec.soil\_evap\_mm & 0.2027\\
{}svd.null[2].vec.canopy\_evap\_mm & 1\\
{}svd.null[2].vec.drainage\_1m\_mm & -0.0192\\
{}balance (W-ET-D) & 47.83, 5.435, 81.74\\
{}ET partition (ET-T-Es-Ec) & 190, 129, 334.7\\
{}clrpca & 0.8858, 0.07317, 0.03278, 0.008299, 0\\
{}clr\_load.transpiration\_mm.0 & -0.218\\
{}clr\_load.transpiration\_mm.1 & 0.035\\
{}clr\_load.transpiration\_mm.2 & -0.22\\
{}clr\_load.soil\_evap\_mm.0 & -0.334\\
{}clr\_load.soil\_evap\_mm.1 & -0.557\\
{}clr\_load.soil\_evap\_mm.2 & -0.449\\
{}clr\_load.canopy\_evap\_mm.0 & -0.034\\
{}clr\_load.canopy\_evap\_mm.1 & 0.801\\
{}clr\_load.canopy\_evap\_mm.2 & -0.21\\
{}clr\_load.R.0 & -0.285\\
{}clr\_load.R.1 & -0.079\\
{}clr\_load.R.2 & 0.839\\
{}clr\_load.drainage\_1m\_mm.0 & 0.871\\
{}clr\_load.drainage\_1m\_mm.1 & -0.2\\
{}clr\_load.drainage\_1m\_mm.2 & 0.04\\
\bottomrule\end{longtable}}

\subsection*{p2\_dpf.json}
\noindent\small Source: \nolinkurl{Manuscripts/results/p2_dpf.json}; 38 values. Workbook sheet \texttt{R\_p2\_dpf}.\par
{\scriptsize\begin{longtable}{p{\dimexpr 0.4557\linewidth-2\tabcolsep\relax}p{\dimexpr 0.5443\linewidth-2\tabcolsep\relax}}\toprule\textbf{quantity} & \textbf{value}\\\midrule\endhead
{}exogenous.terms & lnW/\_TAW, lnalpha, Ifrac, dT\\
{}exogenous.hist[0] & lnW/\_TAW, 0.1926\\
{}exogenous.hist[1] & lnalpha, 0.2601\\
{}exogenous.hist[2] & Ifrac, 0.422\\
{}exogenous.hist[3] & dT, 0.4334\\
{}exogenous.loeo & 0.4334\\
{}exogenous.loso & 0.2171\\
{}exogenous.r2 & 0.5752\\
{}exogenous.coef.const & -12.1\\
{}exogenous.coef.lnW\_TAW & 4.411\\
{}exogenous.coef.lnalpha & -1.214\\
{}exogenous.coef.Ifrac & -5.565\\
{}exogenous.coef.dT & -0.9978\\
{}exogenous.se.const & 0.1018\\
{}exogenous.se.lnW\_TAW & 0.03957\\
{}exogenous.se.lnalpha & 0.0273\\
{}exogenous.se.Ifrac & 0.1049\\
{}exogenous.se.dT & 0.03043\\
{}exogenous.gbm\_loeo & 0.4975\\
{}with\_partition.terms & lnW/\_TAW, lnalpha, Ifrac, dT\\
{}with\_partition.hist[0] & lnW/\_TAW, 0.1926\\
{}with\_partition.hist[1] & lnalpha, 0.2601\\
{}with\_partition.hist[2] & Ifrac, 0.422\\
{}with\_partition.hist[3] & dT, 0.4334\\
{}with\_partition.loeo & 0.4334\\
{}with\_partition.loso & 0.2171\\
{}with\_partition.r2 & 0.5752\\
{}with\_partition.coef.const & -12.1\\
{}with\_partition.coef.lnW\_TAW & 4.411\\
{}with\_partition.coef.lnalpha & -1.214\\
{}with\_partition.coef.Ifrac & -5.565\\
{}with\_partition.coef.dT & -0.9978\\
{}with\_partition.se.const & 0.1018\\
{}with\_partition.se.lnW\_TAW & 0.03957\\
{}with\_partition.se.lnalpha & 0.0273\\
{}with\_partition.se.Ifrac & 0.1049\\
{}with\_partition.se.dT & 0.03043\\
{}with\_partition.gbm\_loeo & 0.5906\\
\bottomrule\end{longtable}}

\subsection*{p2\_et.json}
\noindent\small Source: \nolinkurl{Manuscripts/results/p2_et.json}; 90 values. Workbook sheet \texttt{R\_p2\_et}.\par
{\scriptsize\begin{longtable}{p{\dimexpr 0.4557\linewidth-2\tabcolsep\relax}p{\dimexpr 0.5443\linewidth-2\tabcolsep\relax}}\toprule\textbf{quantity} & \textbf{value}\\\midrule\endhead
{}lnET.terms & lnW, n, lnI1\\
{}lnET.hist[0] & lnW, 0.8008\\
{}lnET.hist[1] & n, 0.8947\\
{}lnET.hist[2] & lnI1, 0.8982\\
{}lnET.loeo & 0.8982\\
{}lnET.loso & 0.9037\\
{}lnET.r2 & 0.9082\\
{}lnET.coef.const & 2.503\\
{}lnET.coef.lnW & 0.6458\\
{}lnET.coef.n & -0.2992\\
{}lnET.coef.lnI1 & 0.02932\\
{}lnET.se.const & 0.02691\\
{}lnET.se.lnW & 0.005084\\
{}lnET.se.n & 0.004524\\
{}lnET.se.lnI1 & 0.001125\\
{}lnET.gbm\_loeo & 0.8789\\
{}lnET.gbm\_loso & 0.8936\\
{}lnET.rmse\_loeo & 0.143\\
{}lnT.terms & lnW, n\\
{}lnT.hist[0] & lnW, 0.5805\\
{}lnT.hist[1] & n, 0.6409\\
{}lnT.loeo & 0.6409\\
{}lnT.loso & 0.6524\\
{}lnT.r2 & 0.6712\\
{}lnT.coef.const & -1.329\\
{}lnT.coef.lnW & 1.127\\
{}lnT.coef.n & -0.5007\\
{}lnT.se.const & 0.06105\\
{}lnT.se.lnW & 0.00874\\
{}lnT.se.n & 0.01243\\
{}lnT.gbm\_loeo & 0.6415\\
{}lnT.gbm\_loso & 0.6592\\
{}lnT.rmse\_loeo & 0.4588\\
{}lnEs.terms & n, mulch, lnrain\\
{}lnEs.hist[0] & n, 0.07031\\
{}lnEs.hist[1] & mulch, 0.2369\\
{}lnEs.hist[2] & lnrain, 0.3084\\
{}lnEs.loeo & 0.3084\\
{}lnEs.loso & 0.3484\\
{}lnEs.r2 & 0.4363\\
{}lnEs.coef.const & 3.271\\
{}lnEs.coef.n & -0.7698\\
{}lnEs.coef.mulch & -0.3804\\
{}lnEs.coef.lnrain & 0.4616\\
{}lnEs.se.const & 0.06683\\
{}lnEs.se.n & 0.01217\\
{}lnEs.se.mulch & 0.006953\\
{}lnEs.se.lnrain & 0.01089\\
{}lnEs.gbm\_loeo & 0.5468\\
{}lnEs.gbm\_loso & 0.4758\\
{}lnEs.rmse\_loeo & 0.4803\\
{}lnEc.terms & lnW, lnrain, Ifrac\\
{}lnEc.hist[0] & lnW, 0.6408\\
{}lnEc.hist[1] & lnrain, 0.6513\\
{}lnEc.hist[2] & Ifrac, 0.6677\\
{}lnEc.loeo & 0.6677\\
{}lnEc.loso & 0.6695\\
{}lnEc.r2 & 0.7035\\
{}lnEc.coef.const & -13.32\\
{}lnEc.coef.lnW & -0.4701\\
{}lnEc.coef.lnrain & 3.157\\
{}lnEc.coef.Ifrac & 4.24\\
{}lnEc.se.const & 0.1448\\
{}lnEc.se.lnW & 0.1351\\
{}lnEc.se.lnrain & 0.1263\\
{}lnEc.se.Ifrac & 0.2092\\
{}lnEc.gbm\_loeo & 0.6914\\
{}lnEc.gbm\_loso & 0.8059\\
{}lnEc.rmse\_loeo & 0.9389\\
{}lnY.terms & lnW, lnI1, lnrain, mulch\\
{}lnY.hist[0] & lnW, 0.5314\\
{}lnY.hist[1] & lnI1, 0.5492\\
{}lnY.hist[2] & lnrain, 0.5992\\
{}lnY.hist[3] & mulch, 0.604\\
{}lnY.loeo & 0.604\\
{}lnY.loso & 0.6169\\
{}lnY.r2 & 0.6354\\
{}lnY.coef.const & -9.681\\
{}lnY.coef.lnW & -0.162\\
{}lnY.coef.lnI1 & 0.3226\\
{}lnY.coef.lnrain & 1.716\\
{}lnY.coef.mulch & 0.1489\\
{}lnY.se.const & 0.1806\\
{}lnY.se.lnW & 0.05911\\
{}lnY.se.lnI1 & 0.008161\\
{}lnY.se.lnrain & 0.04756\\
{}lnY.se.mulch & 0.01476\\
{}lnY.gbm\_loeo & 0.6501\\
{}lnY.gbm\_loso & 0.7425\\
{}lnY.rmse\_loeo & 0.9531\\
\bottomrule\end{longtable}}

\subsection*{p2\_etlaw.json}
\noindent\small Source: \nolinkurl{Manuscripts/results/p2_etlaw.json}; 31 values. Workbook sheet \texttt{R\_p2\_etlaw}.\par
{\scriptsize\begin{longtable}{p{\dimexpr 0.4557\linewidth-2\tabcolsep\relax}p{\dimexpr 0.5443\linewidth-2\tabcolsep\relax}}\toprule\textbf{quantity} & \textbf{value}\\\midrule\endhead
{}sparse\_loeo & 0.8982\\
{}sparse\_loso & 0.9037\\
{}gbm\_loeo & 0.876\\
{}gbm\_loso & 0.8812\\
{}coef\_mean & 2.508, 0.6448, -0.2984, 0.02945\\
{}coef\_sd & 0.0771, 0.01676, 0.01733, 0.002268\\
{}coefS\_mean & 2.498, 0.6465, -0.2993, 0.02912\\
{}coefS\_sd & 0.07797, 0.0124, 0.007435, 0.002257\\
{}rmse\_mm\_loeo & 68.82\\
{}mape\_loeo & 9.179\\
{}by\_exp[0] & E0/\_Baselines, 10, 5.834, -25.17\\
{}by\_exp[1] & E10/\_Amount/\_Timing/\_Soil, 954, 3.256, 1.285\\
{}by\_exp[2] & E11/\_Temp/\_Evap/\_Rain, 1000, 8.999, -62.58\\
{}by\_exp[3] & E12/\_Sensor/\_Trigger/\_Rain, 1000, 3.268, -15.24\\
{}by\_exp[4] & E13/\_N/\_Water/\_Soil, 975, 3.532, -3.65\\
{}by\_exp[5] & E14/\_Mulch/\_Rain, 1000, 4.024, 7.484\\
{}by\_exp[6] & E15/\_Method/\_Depth/\_Amount, 1000, 4.978, 28.96\\
{}by\_exp[7] & E16/\_Edge/\_Cases, 740, 16.6, 34.49\\
{}by\_exp[8] & E1/\_Doses, 189, 3.76, 9.833\\
{}by\_exp[9] & E2/\_TriggerPressure, 197, 3.002, -6.6\\
{}by\_exp[10] & E3/\_SensorDepth, 200, 2.648, -8.047\\
{}by\_exp[11] & E4/\_Methods, 96, 3.557, 9.425\\
{}by\_exp[12] & E5/\_SHP/\_Dynamics, 300, 11.7, -63.8\\
{}by\_exp[13] & E6/\_Mediterranean/\_Pollution, 193, 8.457, 41.52\\
{}by\_exp[14] & E7/\_Mulching, 200, 4.66, 10.62\\
{}by\_exp[15] & E8/\_Rain/\_Intensity/\_Moisture, 1000, 16.72, -11.34\\
{}by\_exp[16] & E9/\_Soil/\_Moisture/\_Rain, 987, 27.4, 12.89\\
{}closureD.r2 & 0.6877\\
{}closureD.rmse & 63.71\\
{}closureD.n & 8054\\
{}closureD.negshare & 0.04817\\
\bottomrule\end{longtable}}

\subsection*{p2\_laws.json}
\noindent\small Source: \nolinkurl{Manuscripts/results/p2_laws.json}; 12 values. Workbook sheet \texttt{R\_p2\_laws}.\par
{\scriptsize\begin{longtable}{p{\dimexpr 0.4557\linewidth-2\tabcolsep\relax}p{\dimexpr 0.5443\linewidth-2\tabcolsep\relax}}\toprule\textbf{quantity} & \textbf{value}\\\midrule\endhead
{}svd4.vec & 1.037, 0.9575, -1, -1.036\\
{}svd4.lo & 1.031, 0.9396, -1, -1.052\\
{}svd4.hi & 1.045, 0.9738, -1, -1.023\\
{}svd4.sv & 176.6, 75.65, 56.89, 5.92\\
{}svd4\_perm\_minSV & 33.33\\
{}svd4\_noE89.vec & 1.021, 0.9688, -1, -1.019\\
{}svd4\_noE89.sv & 158.9, 67.04, 49.48, 3.007\\
{}closure.all\_median\_abs & 0.009134\\
{}closure.p95 & 0.3515\\
{}closure.noE89\_median & 0.006472\\
{}closure.noE89\_p95 & 0.02268\\
{}closure.E89\_median & 0.1643\\
\bottomrule\end{longtable}}

\subsection*{p2\_laws2.json}
\noindent\small Source: \nolinkurl{Manuscripts/results/p2_laws2.json}; 81 values. Workbook sheet \texttt{R\_p2\_laws2}.\par
{\scriptsize\begin{longtable}{p{\dimexpr 0.4557\linewidth-2\tabcolsep\relax}p{\dimexpr 0.5443\linewidth-2\tabcolsep\relax}}\toprule\textbf{quantity} & \textbf{value}\\\midrule\endhead
{}comp\_lattice[0].exp & E11\_Temp\_Evap\_Rain\\
{}comp\_lattice[0].n & 1000\\
{}comp\_lattice[0].irrigation\_mm\_slope & -0.4193\\
{}comp\_lattice[0].irrigation\_mm\_se & 0.01881\\
{}comp\_lattice[0].ET\_mm\_slope & 0.4302\\
{}comp\_lattice[0].ET\_mm\_se & 0.02525\\
{}comp\_lattice[0].drainage\_1m\_mm\_slope & 0.1333\\
{}comp\_lattice[0].drainage\_1m\_mm\_se & 0.006955\\
{}comp\_lattice[0].W\_slope & 0.5807\\
{}comp\_lattice[0].W\_se & 0.01881\\
{}comp\_lattice[1].exp & E12\_Sensor\_Trigger\_Rain\\
{}comp\_lattice[1].n & 1000\\
{}comp\_lattice[1].irrigation\_mm\_slope & -0.673\\
{}comp\_lattice[1].irrigation\_mm\_se & 0.01357\\
{}comp\_lattice[1].ET\_mm\_slope & 0.07769\\
{}comp\_lattice[1].ET\_mm\_se & 0.009059\\
{}comp\_lattice[1].drainage\_1m\_mm\_slope & 0.2316\\
{}comp\_lattice[1].drainage\_1m\_mm\_se & 0.006128\\
{}comp\_lattice[1].W\_slope & 0.327\\
{}comp\_lattice[1].W\_se & 0.01357\\
{}comp\_lattice[2].exp & E14\_Mulch\_Rain\\
{}comp\_lattice[2].n & 1000\\
{}comp\_lattice[2].irrigation\_mm\_slope & -0.6656\\
{}comp\_lattice[2].irrigation\_mm\_se & 0.006611\\
{}comp\_lattice[2].ET\_mm\_slope & -0.01901\\
{}comp\_lattice[2].ET\_mm\_se & 0.009262\\
{}comp\_lattice[2].drainage\_1m\_mm\_slope & 0.3204\\
{}comp\_lattice[2].drainage\_1m\_mm\_se & 0.005513\\
{}comp\_lattice[2].W\_slope & 0.3344\\
{}comp\_lattice[2].W\_se & 0.006611\\
{}comp16.S01.kappa & 0.3033\\
{}comp16.S01.dET & 0.5118\\
{}comp16.S01.dD & 0.186\\
{}comp16.S01.n & 32\\
{}comp16.S02.kappa & 0.1641\\
{}comp16.S02.dET & 0.346\\
{}comp16.S02.dD & 0.4954\\
{}comp16.S02.n & 8\\
{}comp16.S03.kappa & -0.0141\\
{}comp16.S03.dET & 0.4732\\
{}comp16.S03.dD & 0.5502\\
{}comp16.S03.n & 32\\
{}comp16.S04.kappa & 0.08366\\
{}comp16.S04.dET & 0.4733\\
{}comp16.S04.dD & 0.4532\\
{}comp16.S04.n & 32\\
{}comp16.S05.kappa & 0.1021\\
{}comp16.S05.dET & 0.5515\\
{}comp16.S05.dD & 0.3576\\
{}comp16.S05.n & 32\\
{}comp16.S06.kappa & 0.1082\\
{}comp16.S06.dET & 0.6193\\
{}comp16.S06.dD & 0.2813\\
{}comp16.S06.n & 32\\
{}comp16.S07.kappa & 0.2158\\
{}comp16.S07.dET & 0.5423\\
{}comp16.S07.dD & 0.2464\\
{}comp16.S07.n & 32\\
{}comp16.S08.kappa & 0.08991\\
{}comp16.S08.dET & 0.6131\\
{}comp16.S08.dD & 0.3076\\
{}comp16.S08.n & 32\\
{}comp16.S09.kappa & 0.2382\\
{}comp16.S09.dET & 0.5154\\
{}comp16.S09.dD & 0.2529\\
{}comp16.S09.n & 32\\
{}comp16.S10.kappa & 0.3016\\
{}comp16.S10.dET & 0.4768\\
{}comp16.S10.dD & 0.2369\\
{}comp16.S10.n & 32\\
{}kappa16\_pool.est & 0.1589\\
{}kappa16\_pool.lo & 0.1336\\
{}kappa16\_pool.hi & 0.1813\\
{}kap\_by.fEp.0.3 & 0.1927\\
{}kap\_by.fEp.1.2 & 0.1232\\
{}kap\_by.fT.2.0 & 0.188\\
{}kap\_by.fT.6.0 & 0.1289\\
{}kap\_by.fh.50.0 & 0.211\\
{}kap\_by.fh.2000.0 & 0.1067\\
{}kap\_by.fD.1.2 & 0.0791\\
{}kap\_by.fD.30.0 & 0.2408\\
\bottomrule\end{longtable}}

\subsection*{p2\_nested.json}
\noindent\small Source: \nolinkurl{Manuscripts/results/p2_nested.json}; 18 values. Workbook sheet \texttt{R\_p2\_nested}.\par
{\scriptsize\begin{longtable}{p{\dimexpr 0.4557\linewidth-2\tabcolsep\relax}p{\dimexpr 0.5443\linewidth-2\tabcolsep\relax}}\toprule\textbf{quantity} & \textbf{value}\\\midrule\endhead
{}r2 & 0.8946\\
{}sets[0] & lnW, n, lnI1\\
{}sets[1] & lnW, n, lnI1\\
{}sets[2] & lnW, n, lnI1\\
{}sets[3] & lnW, n, lnI1\\
{}sets[4] & lnW, n, lnI1\\
{}sets[5] & lnW, n, lnI1\\
{}sets[6] & lnW, n, lnI1\\
{}sets[7] & lnW, n\\
{}sets[8] & lnW, n, lnI1\\
{}sets[9] & lnW, n, lnI1\\
{}sets[10] & lnW, n, lnI1\\
{}sets[11] & lnW, n, lnI1\\
{}sets[12] & lnW, n, lnI1\\
{}sets[13] & lnW, n\\
{}sets[14] & lnW, n, lnI1\\
{}sets[15] & lnW, n, lnI1\\
{}sets[16] & lnW, n, lnI1\\
\bottomrule\end{longtable}}

\subsection*{p2\_pca.json}
\noindent\small Source: \nolinkurl{Manuscripts/results/p2_pca.json}; 4 values. Workbook sheet \texttt{R\_p2\_pca}.\par
{\scriptsize\begin{longtable}{p{\dimexpr 0.4557\linewidth-2\tabcolsep\relax}p{\dimexpr 0.5443\linewidth-2\tabcolsep\relax}}\toprule\textbf{quantity} & \textbf{value}\\\midrule\endhead
{}evr & 0.8858, 0.07317, 0.03278\\
{}load[0] & -0.2184, -0.3338, -0.03369, -0.2851, 0.8709\\
{}load[1] & 0.03458, -0.5574, 0.8013, -0.07872, -0.1997\\
{}load[2] & -0.2197, -0.4488, -0.2104, 0.8393, 0.03956\\
\bottomrule\end{longtable}}

\subsection*{p2\_pca\_sens.json}
\noindent\small Source: \nolinkurl{Manuscripts/results/p2_pca_sens.json}; 10 values. Workbook sheet \texttt{R\_p2\_pca\_sens}.\par
{\scriptsize\begin{longtable}{p{\dimexpr 0.4557\linewidth-2\tabcolsep\relax}p{\dimexpr 0.5443\linewidth-2\tabcolsep\relax}}\toprule\textbf{quantity} & \textbf{value}\\\midrule\endhead
{}sens.0.001 & 0.8858, 0.07317, 0.03278, 0.008299\\
{}sens.0.1 & 0.7435, 0.1732, 0.06629, 0.01701\\
{}sens.1 & 0.6274, 0.2481, 0.09837, 0.02614\\
{}sens.5 & 0.565, 0.2511, 0.1481, 0.03581\\
{}sens.10 & 0.5233, 0.271, 0.168, 0.03776\\
{}clip1\_evr & 0.6274, 0.2481, 0.09837\\
{}clip1\_load[0] & -0.1982, -0.3672, 0.06212, -0.338, 0.8413\\
{}clip1\_load[1] & 0.08674, -0.5565, 0.7747, -0.01804, -0.2869\\
{}clip1\_load[2] & -0.2647, -0.4327, -0.2259, 0.826, 0.09732\\
{}cor\_PC1 & 0.4828, 0.454, 0.9791\\
\bottomrule\end{longtable}}

\subsection*{p2\_robust.json}
\noindent\small Source: \nolinkurl{Manuscripts/results/p2_robust.json}; 121 values. Workbook sheet \texttt{R\_p2\_robust}.\par
{\scriptsize\begin{longtable}{p{\dimexpr 0.4557\linewidth-2\tabcolsep\relax}p{\dimexpr 0.5443\linewidth-2\tabcolsep\relax}}\toprule\textbf{quantity} & \textbf{value}\\\midrule\endhead
{}terms & lnW, n, lnI1\\
{}bootstrap.n\_resamples & 2000\\
{}bootstrap.law\_loeo & 0.8982\\
{}bootstrap.gbm3\_loeo & 0.8662\\
{}bootstrap.diff & 0.03198\\
{}bootstrap.diff\_mean & 0.03612\\
{}bootstrap.ci95 & 0.004383, 0.0867\\
{}without\_E8\_E9.n\_runs & 8054\\
{}without\_E8\_E9.n\_experiments & 15\\
{}without\_E8\_E9.loeo & 0.8818\\
{}without\_E8\_E9.loso & 0.8811\\
{}without\_E8\_E9.rmse\_mm & 65.88\\
{}without\_E8\_E9.mape\_pct & 5.787\\
{}without\_E8\_E9.steps[0].terms & lnW\\
{}without\_E8\_E9.steps[0].loeo & 0.5998\\
{}without\_E8\_E9.steps[0].rmse\_mm & 137.9\\
{}without\_E8\_E9.steps[1].terms & lnW, n\\
{}without\_E8\_E9.steps[1].loeo & 0.8567\\
{}without\_E8\_E9.steps[1].rmse\_mm & 73.84\\
{}without\_E8\_E9.steps[2].terms & lnW, n, lnI1\\
{}without\_E8\_E9.steps[2].loeo & 0.8818\\
{}without\_E8\_E9.steps[2].rmse\_mm & 65.88\\
{}without\_E8\_E9.gbm3\_loeo & 0.7781\\
{}table1.rain\_mm.mean & 417\\
{}table1.rain\_mm.sd & 143.3\\
{}table1.rain\_mm.q25 & 297.4\\
{}table1.rain\_mm.median & 464.4\\
{}table1.rain\_mm.q75 & 464.4\\
{}table1.irrigation\_mm.mean & 179.7\\
{}table1.irrigation\_mm.sd & 191.4\\
{}table1.irrigation\_mm.q25 & 0\\
{}table1.irrigation\_mm.median & 178.2\\
{}table1.irrigation\_mm.q75 & 268.4\\
{}table1.W.mean & 596.7\\
{}table1.W.sd & 251\\
{}table1.W.q25 & 464.4\\
{}table1.W.median & 654.5\\
{}table1.W.q75 & 722.1\\
{}table1.ET\_mm.mean & 546.1\\
{}table1.ET\_mm.sd & 195.7\\
{}table1.ET\_mm.q25 & 385.9\\
{}table1.ET\_mm.median & 611.5\\
{}table1.ET\_mm.q75 & 670.7\\
{}table1.drainage\_1m\_mm.mean & 60.09\\
{}table1.drainage\_1m\_mm.sd & 105.8\\
{}table1.drainage\_1m\_mm.q25 & 2.803\\
{}table1.drainage\_1m\_mm.median & 33.51\\
{}table1.drainage\_1m\_mm.q75 & 77.62\\
{}table1.R.mean & 163.7\\
{}table1.R.sd & 96.47\\
{}table1.R.q25 & 102.8\\
{}table1.R.median & 129\\
{}table1.R.q75 & 222.5\\
{}R\_by\_experiment.E0\_Baselines.runs & 10\\
{}R\_by\_experiment.E0\_Baselines.R\_median & 106.9\\
{}R\_by\_experiment.E0\_Baselines.R\_ET\_median & 0.2398\\
{}R\_by\_experiment.E0\_Baselines.dS\_W\_median & -0.02058\\
{}R\_by\_experiment.E10\_Amount\_Timing\_Soil.runs & 954\\
{}R\_by\_experiment.E10\_Amount\_Timing\_Soil.R\_median & 118.3\\
{}R\_by\_experiment.E10\_Amount\_Timing\_Soil.R\_ET\_median & 0.2295\\
{}R\_by\_experiment.E10\_Amount\_Timing\_Soil.dS\_W\_median & -0.01078\\
{}R\_by\_experiment.E11\_Temp\_Evap\_Rain.runs & 1000\\
{}R\_by\_experiment.E11\_Temp\_Evap\_Rain.R\_median & 296.3\\
{}R\_by\_experiment.E11\_Temp\_Evap\_Rain.R\_ET\_median & 0.3751\\
{}R\_by\_experiment.E11\_Temp\_Evap\_Rain.dS\_W\_median & -0.009339\\
{}R\_by\_experiment.E12\_Sensor\_Trigger\_Rain.runs & 1000\\
{}R\_by\_experiment.E12\_Sensor\_Trigger\_Rain.R\_median & 217.6\\
{}R\_by\_experiment.E12\_Sensor\_Trigger\_Rain.R\_ET\_median & 0.3186\\
{}R\_by\_experiment.E12\_Sensor\_Trigger\_Rain.dS\_W\_median & -0.008244\\
{}R\_by\_experiment.E13\_N\_Water\_Soil.runs & 975\\
{}R\_by\_experiment.E13\_N\_Water\_Soil.R\_median & 124.5\\
{}R\_by\_experiment.E13\_N\_Water\_Soil.R\_ET\_median & 0.284\\
{}R\_by\_experiment.E13\_N\_Water\_Soil.dS\_W\_median & -0.006363\\
{}R\_by\_experiment.E14\_Mulch\_Rain.runs & 1000\\
{}R\_by\_experiment.E14\_Mulch\_Rain.R\_median & 262.1\\
{}R\_by\_experiment.E14\_Mulch\_Rain.R\_ET\_median & 0.409\\
{}R\_by\_experiment.E14\_Mulch\_Rain.dS\_W\_median & 0.0008958\\
{}R\_by\_experiment.E15\_Method\_Depth\_Amount.runs & 1000\\
{}R\_by\_experiment.E15\_Method\_Depth\_Amount.R\_median & 103.2\\
{}R\_by\_experiment.E15\_Method\_Depth\_Amount.R\_ET\_median & 0.1689\\
{}R\_by\_experiment.E15\_Method\_Depth\_Amount.dS\_W\_median & 0.003355\\
{}R\_by\_experiment.E16\_Edge\_Cases.runs & 740\\
{}R\_by\_experiment.E16\_Edge\_Cases.R\_median & 157\\
{}R\_by\_experiment.E16\_Edge\_Cases.R\_ET\_median & 0.2964\\
{}R\_by\_experiment.E16\_Edge\_Cases.dS\_W\_median & -0.01008\\
{}R\_by\_experiment.E1\_Doses.runs & 189\\
{}R\_by\_experiment.E1\_Doses.R\_median & 117.6\\
{}R\_by\_experiment.E1\_Doses.R\_ET\_median & 0.2123\\
{}R\_by\_experiment.E1\_Doses.dS\_W\_median & -0.001241\\
{}R\_by\_experiment.E2\_TriggerPressure.runs & 197\\
{}R\_by\_experiment.E2\_TriggerPressure.R\_median & 167.5\\
{}R\_by\_experiment.E2\_TriggerPressure.R\_ET\_median & 0.2979\\
{}R\_by\_experiment.E2\_TriggerPressure.dS\_W\_median & -0.008657\\
{}R\_by\_experiment.E3\_SensorDepth.runs & 200\\
{}R\_by\_experiment.E3\_SensorDepth.R\_median & 201.8\\
{}R\_by\_experiment.E3\_SensorDepth.R\_ET\_median & 0.3213\\
{}R\_by\_experiment.E3\_SensorDepth.dS\_W\_median & -0.008657\\
{}R\_by\_experiment.E4\_Methods.runs & 96\\
{}R\_by\_experiment.E4\_Methods.R\_median & 103.4\\
{}R\_by\_experiment.E4\_Methods.R\_ET\_median & 0.1863\\
{}R\_by\_experiment.E4\_Methods.dS\_W\_median & -0.003341\\
{}R\_by\_experiment.E5\_SHP\_Dynamics.runs & 300\\
{}R\_by\_experiment.E5\_SHP\_Dynamics.R\_median & 161.5\\
{}R\_by\_experiment.E5\_SHP\_Dynamics.R\_ET\_median & 0.2882\\
{}R\_by\_experiment.E5\_SHP\_Dynamics.dS\_W\_median & -0.007413\\
{}R\_by\_experiment.E6\_Mediterranean\_Pollution.runs & 193\\
{}R\_by\_experiment.E6\_Mediterranean\_Pollution.R\_median & 70.43\\
{}R\_by\_experiment.E6\_Mediterranean\_Pollution.R\_ET\_median & 0.1435\\
{}R\_by\_experiment.E6\_Mediterranean\_Pollution.dS\_W\_median & -0.01191\\
{}R\_by\_experiment.E7\_Mulching.runs & 200\\
{}R\_by\_experiment.E7\_Mulching.R\_median & 226.3\\
{}R\_by\_experiment.E7\_Mulching.R\_ET\_median & 0.4091\\
{}R\_by\_experiment.E7\_Mulching.dS\_W\_median & -0.005667\\
{}R\_by\_experiment.E8\_Rain\_Intensity\_Moisture.runs & 1000\\
{}R\_by\_experiment.E8\_Rain\_Intensity\_Moisture.R\_median & 39.29\\
{}R\_by\_experiment.E8\_Rain\_Intensity\_Moisture.R\_ET\_median & 0.1438\\
{}R\_by\_experiment.E8\_Rain\_Intensity\_Moisture.dS\_W\_median & -0.1335\\
{}R\_by\_experiment.E9\_Soil\_Moisture\_Rain.runs & 987\\
{}R\_by\_experiment.E9\_Soil\_Moisture\_Rain.R\_median & 81.23\\
{}R\_by\_experiment.E9\_Soil\_Moisture\_Rain.R\_ET\_median & 0.3736\\
{}R\_by\_experiment.E9\_Soil\_Moisture\_Rain.dS\_W\_median & -0.01399\\
\bottomrule\end{longtable}}

\subsection*{p2\_s02\_excl.json}
\noindent\small Source: \nolinkurl{Manuscripts/results/p2_s02_excl.json}; 27 values. Workbook sheet \texttt{R\_p2\_s02\_excl}.\par
{\scriptsize\begin{longtable}{p{\dimexpr 0.4557\linewidth-2\tabcolsep\relax}p{\dimexpr 0.5443\linewidth-2\tabcolsep\relax}}\toprule\textbf{quantity} & \textbf{value}\\\midrule\endhead
{}terms & lnW, n, lnI1\\
{}full.loeo & 0.8982\\
{}full.loso & 0.9037\\
{}full.coef.const & 2.503\\
{}full.coef.lnW & 0.6458\\
{}full.coef.n & -0.2992\\
{}full.coef.lnI1 & 0.02932\\
{}full.R\_mean & 163.7\\
{}full.R\_share\_of\_ET & 0.2998\\
{}without\_S02.n\_runs & 9651\\
{}without\_S02.loeo & 0.8972\\
{}without\_S02.loso & 0.9037\\
{}without\_S02.coef.const & 2.461\\
{}without\_S02.coef.lnW & 0.6519\\
{}without\_S02.coef.n & -0.2972\\
{}without\_S02.coef.lnI1 & 0.02947\\
{}without\_S02.R\_mean & 163.1\\
{}without\_S02.R\_share\_of\_ET & 0.2988\\
{}without\_S02\_S03.n\_runs & 9145\\
{}without\_S02\_S03.loeo & 0.8931\\
{}without\_S02\_S03.loso & 0.8976\\
{}without\_S02\_S03.coef.const & 2.548\\
{}without\_S02\_S03.coef.lnW & 0.6419\\
{}without\_S02\_S03.coef.n & -0.3168\\
{}without\_S02\_S03.coef.lnI1 & 0.03039\\
{}without\_S02\_S03.R\_mean & 165.4\\
{}without\_S02\_S03.R\_share\_of\_ET & 0.2958\\
\bottomrule\end{longtable}}


\section{Where every table and figure comes from}
{\scriptsize\begin{longtable}{p{\dimexpr 0.0955\linewidth-2\tabcolsep\relax}p{\dimexpr 0.2293\linewidth-2\tabcolsep\relax}p{\dimexpr 0.3185\linewidth-2\tabcolsep\relax}p{\dimexpr 0.3567\linewidth-2\tabcolsep\relax}}\toprule\textbf{Item} & \textbf{How it is made} & \textbf{Script(s)} & \textbf{Input and output}\\\midrule\endhead
{}Table 1 & Descriptive table (definitions, design, parameters or tests); no computation & - & \\
{}Figure 1 & Figure drawn by the script from the input data & Manuscripts/\allowbreak{}reproducibility/\allowbreak{}make\_\allowbreak{}a4.py;  Manuscripts/\allowbreak{}reproducibility/\allowbreak{}analysis/\allowbreak{}p2\_\allowbreak{}figs.py;  Manuscripts/\allowbreak{}reproducibility/\allowbreak{}analysis/\allowbreak{}p2\_\allowbreak{}fig1\_\allowbreak{}new.py & Data/\allowbreak{}daisy\_\allowbreak{}ensemble.csv.gz (10 190 runs);  Manuscripts/\allowbreak{}figures/\allowbreak{}p2/\allowbreak{}fig01\_\allowbreak{}pipeline\_\allowbreak{}A4.pdf\\
{}Table 2 & Descriptive table (definitions, design, parameters or tests); no computation & - & \\
{}Table 3 & Table typed from the result file (27 of its numbers found there) & Manuscripts/\allowbreak{}reproducibility/\allowbreak{}p3\_\allowbreak{}js/\allowbreak{}run\_\allowbreak{}drip.js;  Manuscripts/\allowbreak{}reproducibility/\allowbreak{}analysis/\allowbreak{}p3\_\allowbreak{}lit\_\allowbreak{}figs.py;  Manuscripts/\allowbreak{}reproducibility/\allowbreak{}analysis/\allowbreak{}make\_\allowbreak{}supplement\_\allowbreak{}data.py & Data/\allowbreak{}daisy\_\allowbreak{}ensemble.csv.gz (10 190 runs);  Manuscripts/\allowbreak{}results/\allowbreak{}drip\_\allowbreak{}out.json\\
{}Figure 2 & Figure drawn by the script from the input data & Manuscripts/\allowbreak{}reproducibility/\allowbreak{}make\_\allowbreak{}a4.py;  Manuscripts/\allowbreak{}reproducibility/\allowbreak{}analysis/\allowbreak{}p1\_\allowbreak{}figs\_\allowbreak{}d.py;  Manuscripts/\allowbreak{}reproducibility/\allowbreak{}analysis/\allowbreak{}p1\_\allowbreak{}fig1\_\allowbreak{}new.py & Data/\allowbreak{}daisy\_\allowbreak{}ensemble.csv.gz (10 190 runs);  Manuscripts/\allowbreak{}figures/\allowbreak{}p1/\allowbreak{}fig01\_\allowbreak{}framework\_\allowbreak{}A4.pdf\\
{}Figure 3 & Figure drawn by the script from the input data & Manuscripts/\allowbreak{}reproducibility/\allowbreak{}analysis/\allowbreak{}p1\_\allowbreak{}figs\_\allowbreak{}a.py & Data/\allowbreak{}daisy\_\allowbreak{}ensemble.csv.gz (10 190 runs);  Manuscripts/\allowbreak{}figures/\allowbreak{}p1/\allowbreak{}fig02\_\allowbreak{}hydraulics.pdf\\
{}Table 4 & Table written by the script (tex/data\_table.tex) & Manuscripts/\allowbreak{}reproducibility/\allowbreak{}analysis/\allowbreak{}data\_\allowbreak{}overview.py & Data/\allowbreak{}Comprehensive\_\allowbreak{}Summary/\allowbreak{}00\_\allowbreak{}summary\_\allowbreak{}all\_\allowbreak{}simulations.csv.gz\\
{}Figure 4 & Figure drawn by the script from the input data & Manuscripts/\allowbreak{}reproducibility/\allowbreak{}analysis/\allowbreak{}data\_\allowbreak{}overview.py & Data/\allowbreak{}daisy\_\allowbreak{}ensemble.csv.gz (10 190 runs);  Manuscripts/\allowbreak{}figures/\allowbreak{}data/\allowbreak{}fig\_\allowbreak{}data\_\allowbreak{}overview.pdf\\
{}Table 5 & Table typed in the manuscript from the result files of the script & p2\_\allowbreak{}robust.py;  p2\_\allowbreak{}robust.js & Data/\allowbreak{}daisy\_\allowbreak{}ensemble.csv.gz (10 190 runs)\\
{}Table 6 & Table typed from the result file (11 of its numbers found there) & Manuscripts/\allowbreak{}reproducibility/\allowbreak{}audit\_\allowbreak{}numbers.py;  Manuscripts/\allowbreak{}reproducibility/\allowbreak{}p3\_\allowbreak{}js/\allowbreak{}e1\_\allowbreak{}heldout.js;  Manuscripts/\allowbreak{}reproducibility/\allowbreak{}analysis/\allowbreak{}p3\_\allowbreak{}figs.py & Data/\allowbreak{}daisy\_\allowbreak{}ensemble.csv.gz (10 190 runs);  Manuscripts/\allowbreak{}results/\allowbreak{}e12\_\allowbreak{}pred.json\\
{}Figure 5 & Figure drawn by the script from the input data & Manuscripts/\allowbreak{}reproducibility/\allowbreak{}analysis/\allowbreak{}p2\_\allowbreak{}figfail.py & Data/\allowbreak{}daisy\_\allowbreak{}ensemble.csv.gz (10 190 runs);  Manuscripts/\allowbreak{}figures/\allowbreak{}p2/\allowbreak{}figS\_\allowbreak{}failed.pdf\\
{}Table 7 & Descriptive table (definitions, design, parameters or tests); no computation & - & \\
{}Figure 6 & Figure drawn by the script from the input data & Manuscripts/\allowbreak{}reproducibility/\allowbreak{}analysis/\allowbreak{}p1\_\allowbreak{}figs\_\allowbreak{}a.py & Data/\allowbreak{}daisy\_\allowbreak{}ensemble.csv.gz (10 190 runs);  Manuscripts/\allowbreak{}figures/\allowbreak{}p1/\allowbreak{}fig03\_\allowbreak{}memory\_\allowbreak{}grid.pdf\\
{}Figure 7 & Figure drawn by the script from the input data & Manuscripts/\allowbreak{}reproducibility/\allowbreak{}analysis/\allowbreak{}p1\_\allowbreak{}figs\_\allowbreak{}b.py & Data/\allowbreak{}daisy\_\allowbreak{}ensemble.csv.gz (10 190 runs);  Manuscripts/\allowbreak{}figures/\allowbreak{}p1/\allowbreak{}fig04\_\allowbreak{}fates.pdf\\
{}Table 8 & Table typed from the result file (55 of its numbers found there) & Manuscripts/\allowbreak{}reproducibility/\allowbreak{}p3\_\allowbreak{}js/\allowbreak{}run\_\allowbreak{}drip.js;  Manuscripts/\allowbreak{}reproducibility/\allowbreak{}analysis/\allowbreak{}p3\_\allowbreak{}lit\_\allowbreak{}figs.py;  Manuscripts/\allowbreak{}reproducibility/\allowbreak{}analysis/\allowbreak{}make\_\allowbreak{}supplement\_\allowbreak{}data.py;  p1\_\allowbreak{}core.py;  p1\_\allowbreak{}core.js & Data/\allowbreak{}daisy\_\allowbreak{}ensemble.csv.gz (10 190 runs);  Manuscripts/\allowbreak{}results/\allowbreak{}drip\_\allowbreak{}out.json\\
{}Figure 8 & Figure drawn by the script from the input data & Manuscripts/\allowbreak{}reproducibility/\allowbreak{}analysis/\allowbreak{}p1\_\allowbreak{}figs\_\allowbreak{}b.py & Data/\allowbreak{}daisy\_\allowbreak{}ensemble.csv.gz (10 190 runs);  Manuscripts/\allowbreak{}figures/\allowbreak{}p1/\allowbreak{}fig05\_\allowbreak{}predictors.pdf\\
{}Figure 9 & Figure drawn by the script from the input data & Manuscripts/\allowbreak{}reproducibility/\allowbreak{}analysis/\allowbreak{}p1\_\allowbreak{}figs\_\allowbreak{}b.py;  p1\_\allowbreak{}rate\_\allowbreak{}law.py;  p1\_\allowbreak{}rate\_\allowbreak{}law.js & Data/\allowbreak{}daisy\_\allowbreak{}ensemble.csv.gz (10 190 runs);  Manuscripts/\allowbreak{}figures/\allowbreak{}p1/\allowbreak{}fig06\_\allowbreak{}mechanism.pdf\\
{}Table 9 & Table typed from the result file (8 of its numbers found there) & Manuscripts/\allowbreak{}reproducibility/\allowbreak{}p3\_\allowbreak{}js/\allowbreak{}run\_\allowbreak{}drip.js;  Manuscripts/\allowbreak{}reproducibility/\allowbreak{}analysis/\allowbreak{}p3\_\allowbreak{}lit\_\allowbreak{}figs.py;  Manuscripts/\allowbreak{}reproducibility/\allowbreak{}analysis/\allowbreak{}make\_\allowbreak{}supplement\_\allowbreak{}data.py;  p1\_\allowbreak{}rate\_\allowbreak{}law.py;  p1\_\allowbreak{}rate\_\allowbreak{}law.js & Data/\allowbreak{}daisy\_\allowbreak{}ensemble.csv.gz (10 190 runs);  Manuscripts/\allowbreak{}results/\allowbreak{}drip\_\allowbreak{}out.json\\
{}Figure 10 & Figure drawn by the script from the input data & Manuscripts/\allowbreak{}reproducibility/\allowbreak{}analysis/\allowbreak{}p1\_\allowbreak{}figs\_\allowbreak{}c.py & Data/\allowbreak{}daisy\_\allowbreak{}ensemble.csv.gz (10 190 runs);  Manuscripts/\allowbreak{}figures/\allowbreak{}p1/\allowbreak{}fig08\_\allowbreak{}replication.pdf\\
{}Figure 11 & Figure drawn by the script from the input data & Manuscripts/\allowbreak{}reproducibility/\allowbreak{}analysis/\allowbreak{}p1\_\allowbreak{}figs\_\allowbreak{}d.py & Data/\allowbreak{}daisy\_\allowbreak{}ensemble.csv.gz (10 190 runs);  Manuscripts/\allowbreak{}figures/\allowbreak{}p1/\allowbreak{}fig10\_\allowbreak{}robustness.pdf\\
{}Table 10 & Table typed from the result file (14 of its numbers found there) & p1\_\allowbreak{}stats.py;  p1\_\allowbreak{}figs\_\allowbreak{}d.py;  p1\_\allowbreak{}rank.js & Data/\allowbreak{}daisy\_\allowbreak{}ensemble.csv.gz (10 190 runs);  Results/\allowbreak{}07\_\allowbreak{}Derived\_\allowbreak{}Data\_\allowbreak{}Tables/\allowbreak{}pareto\_\allowbreak{}runs.csv\\
{}Figure 12 & Figure drawn by the script from the input data & Manuscripts/\allowbreak{}reproducibility/\allowbreak{}analysis/\allowbreak{}p1\_\allowbreak{}figs\_\allowbreak{}c.py & Data/\allowbreak{}daisy\_\allowbreak{}ensemble.csv.gz (10 190 runs);  Manuscripts/\allowbreak{}figures/\allowbreak{}p1/\allowbreak{}fig07\_\allowbreak{}modifiers.pdf\\
{}Table 11 & Table typed from the result file (8 of its numbers found there) & Manuscripts/\allowbreak{}reproducibility/\allowbreak{}analysis/\allowbreak{}p1\_\allowbreak{}figs\_\allowbreak{}c.py;  p1\_\allowbreak{}modifiers.js & Data/\allowbreak{}daisy\_\allowbreak{}ensemble.csv.gz (10 190 runs);  Manuscripts/\allowbreak{}results/\allowbreak{}p1\_\allowbreak{}modifiers.json\\
{}Figure 13 & Figure drawn by the script from the input data & Manuscripts/\allowbreak{}reproducibility/\allowbreak{}analysis/\allowbreak{}p1\_\allowbreak{}figs\_\allowbreak{}e.py & Data/\allowbreak{}daisy\_\allowbreak{}ensemble.csv.gz (10 190 runs);  Manuscripts/\allowbreak{}figures/\allowbreak{}p1/\allowbreak{}fig11\_\allowbreak{}e11\_\allowbreak{}surface.pdf\\
{}Table 12 & Table typed from the result file (7 of its numbers found there) & Manuscripts/\allowbreak{}reproducibility/\allowbreak{}p3\_\allowbreak{}js/\allowbreak{}run\_\allowbreak{}drip.js;  Manuscripts/\allowbreak{}reproducibility/\allowbreak{}analysis/\allowbreak{}p3\_\allowbreak{}lit\_\allowbreak{}figs.py;  Manuscripts/\allowbreak{}reproducibility/\allowbreak{}analysis/\allowbreak{}make\_\allowbreak{}supplement\_\allowbreak{}data.py & Data/\allowbreak{}daisy\_\allowbreak{}ensemble.csv.gz (10 190 runs);  Manuscripts/\allowbreak{}results/\allowbreak{}drip\_\allowbreak{}out.json\\
{}Figure 14 & Figure drawn by the script from the input data & Manuscripts/\allowbreak{}reproducibility/\allowbreak{}analysis/\allowbreak{}p2\_\allowbreak{}figs.py;  p2\_\allowbreak{}a.py;  p2\_\allowbreak{}laws.py;  p2\_\allowbreak{}a.js;  p2\_\allowbreak{}laws.js & Data/\allowbreak{}daisy\_\allowbreak{}ensemble.csv.gz (10 190 runs);  Manuscripts/\allowbreak{}figures/\allowbreak{}p2/\allowbreak{}fig02\_\allowbreak{}invariants.pdf\\
{}Table 13 & Table typed from the result file (34 of its numbers found there) & Manuscripts/\allowbreak{}reproducibility/\allowbreak{}p3\_\allowbreak{}js/\allowbreak{}run\_\allowbreak{}drip.js;  Manuscripts/\allowbreak{}reproducibility/\allowbreak{}analysis/\allowbreak{}p3\_\allowbreak{}lit\_\allowbreak{}figs.py;  Manuscripts/\allowbreak{}reproducibility/\allowbreak{}analysis/\allowbreak{}make\_\allowbreak{}supplement\_\allowbreak{}data.py & Data/\allowbreak{}daisy\_\allowbreak{}ensemble.csv.gz (10 190 runs);  Manuscripts/\allowbreak{}results/\allowbreak{}drip\_\allowbreak{}out.json\\
{}Figure 15 & Figure drawn by the script from the input data & Manuscripts/\allowbreak{}reproducibility/\allowbreak{}analysis/\allowbreak{}p2\_\allowbreak{}figs.py;  p2\_\allowbreak{}R.py;  p2\_\allowbreak{}R.js & Data/\allowbreak{}daisy\_\allowbreak{}ensemble.csv.gz (10 190 runs);  Manuscripts/\allowbreak{}figures/\allowbreak{}p2/\allowbreak{}fig03\_\allowbreak{}hidden\_\allowbreak{}term.pdf\\
{}Table 14 & Table typed from the result file (22 of its numbers found there) & Manuscripts/\allowbreak{}reproducibility/\allowbreak{}run\_\allowbreak{}all.sh;  Manuscripts/\allowbreak{}reproducibility/\allowbreak{}audit\_\allowbreak{}numbers.py;  Manuscripts/\allowbreak{}reproducibility/\allowbreak{}analysis/\allowbreak{}p1\_\allowbreak{}core.py;  p2\_\allowbreak{}robust.py;  p2\_\allowbreak{}robust.js & Data/\allowbreak{}daisy\_\allowbreak{}ensemble.csv.gz (10 190 runs);  Manuscripts/\allowbreak{}results/\allowbreak{}p1\_\allowbreak{}core.json\\
{}Figure 16 & Figure drawn by the script from the input data & Manuscripts/\allowbreak{}reproducibility/\allowbreak{}analysis/\allowbreak{}p2\_\allowbreak{}fig4.py;  Manuscripts/\allowbreak{}reproducibility/\allowbreak{}analysis/\allowbreak{}p2\_\allowbreak{}figs.py;  p2\_\allowbreak{}a.py;  p2\_\allowbreak{}pca.js;  p2\_\allowbreak{}pca\_\allowbreak{}sens.js & Data/\allowbreak{}daisy\_\allowbreak{}ensemble.csv.gz (10 190 runs);  Manuscripts/\allowbreak{}figures/\allowbreak{}p2/\allowbreak{}fig04\_\allowbreak{}manifold.pdf\\
{}Table 15 & Table typed in the manuscript from the result files of the script & p2\_\allowbreak{}a.py;  p2\_\allowbreak{}figs.py;  p2\_\allowbreak{}fig4.py;  p2\_\allowbreak{}pca.js;  p2\_\allowbreak{}pca\_\allowbreak{}sens.js & Data/\allowbreak{}daisy\_\allowbreak{}ensemble.csv.gz (10 190 runs)\\
{}Figure 17 & Figure drawn by the script from the input data & Manuscripts/\allowbreak{}reproducibility/\allowbreak{}analysis/\allowbreak{}p2\_\allowbreak{}fig7.py & Data/\allowbreak{}daisy\_\allowbreak{}ensemble.csv.gz (10 190 runs);  Manuscripts/\allowbreak{}figures/\allowbreak{}p2/\allowbreak{}fig07\_\allowbreak{}fates\_\allowbreak{}by\_\allowbreak{}soil.pdf\\
{}Figure 18 & Figure drawn by the script from the input data & Manuscripts/\allowbreak{}reproducibility/\allowbreak{}analysis/\allowbreak{}p2\_\allowbreak{}figs2.py & Data/\allowbreak{}daisy\_\allowbreak{}ensemble.csv.gz (10 190 runs);  Manuscripts/\allowbreak{}figures/\allowbreak{}p2/\allowbreak{}fig05\_\allowbreak{}ET\_\allowbreak{}law.pdf\\
{}Table 16 & Table typed in the manuscript from the result files of the script & (see the reproducibility appendix of the paper) & Data/\allowbreak{}daisy\_\allowbreak{}ensemble.csv.gz (10 190 runs)\\
{}Table 17 & Table typed from the result file (18 of its numbers found there) & Manuscripts/\allowbreak{}reproducibility/\allowbreak{}run\_\allowbreak{}all.sh;  Manuscripts/\allowbreak{}reproducibility/\allowbreak{}make\_\allowbreak{}e12.py;  Manuscripts/\allowbreak{}reproducibility/\allowbreak{}audit\_\allowbreak{}numbers.py & Data/\allowbreak{}daisy\_\allowbreak{}ensemble.csv.gz (10 190 runs);  Manuscripts/\allowbreak{}results/\allowbreak{}e12.json\\
{}Table 18 & Table typed in the manuscript from the result files of the script & (see the reproducibility appendix of the paper) & Data/\allowbreak{}daisy\_\allowbreak{}ensemble.csv.gz (10 190 runs)\\
{}Table 19 & Table typed from the result file (12 of its numbers found there) & Manuscripts/\allowbreak{}reproducibility/\allowbreak{}p3\_\allowbreak{}js/\allowbreak{}run\_\allowbreak{}drip.js;  Manuscripts/\allowbreak{}reproducibility/\allowbreak{}analysis/\allowbreak{}p3\_\allowbreak{}lit\_\allowbreak{}figs.py;  Manuscripts/\allowbreak{}reproducibility/\allowbreak{}analysis/\allowbreak{}make\_\allowbreak{}supplement\_\allowbreak{}data.py & Data/\allowbreak{}daisy\_\allowbreak{}ensemble.csv.gz (10 190 runs);  Manuscripts/\allowbreak{}results/\allowbreak{}drip\_\allowbreak{}out.json\\
{}Table 20 & Table typed from the result file (3 of its numbers found there) & Manuscripts/\allowbreak{}reproducibility/\allowbreak{}analysis/\allowbreak{}p4\_\allowbreak{}exact.py;  Manuscripts/\allowbreak{}reproducibility/\allowbreak{}analysis/\allowbreak{}p4\_\allowbreak{}reach.py;  Manuscripts/\allowbreak{}reproducibility/\allowbreak{}analysis/\allowbreak{}make\_\allowbreak{}supplement\_\allowbreak{}data.py & Data/\allowbreak{}daisy\_\allowbreak{}ensemble.csv.gz (10 190 runs);  Manuscripts/\allowbreak{}results/\allowbreak{}p4\_\allowbreak{}tableA1\_\allowbreak{}reach.csv\\
{}Table 21 & Table typed from the result file (23 of its numbers found there) & Manuscripts/\allowbreak{}reproducibility/\allowbreak{}p3\_\allowbreak{}js/\allowbreak{}run\_\allowbreak{}drip.js;  Manuscripts/\allowbreak{}reproducibility/\allowbreak{}analysis/\allowbreak{}p3\_\allowbreak{}lit\_\allowbreak{}figs.py;  Manuscripts/\allowbreak{}reproducibility/\allowbreak{}analysis/\allowbreak{}make\_\allowbreak{}supplement\_\allowbreak{}data.py;  p2\_\allowbreak{}et.py;  p2\_\allowbreak{}dpf.py;  p2\_\allowbreak{}figs2.py;  p2\_\allowbreak{}nested.py;  p2\_\allowbreak{}et.js;  p2\_\allowbreak{}dpf.js;  p2\_\allowbreak{}etlaw.js;  p2\_\allowbreak{}nested.js;  p2\_\allowbreak{}s02\_\allowbreak{}excl.py;  p2\_\allowbreak{}s02\_\allowbreak{}excl.js & Data/\allowbreak{}daisy\_\allowbreak{}ensemble.csv.gz (10 190 runs);  Manuscripts/\allowbreak{}results/\allowbreak{}drip\_\allowbreak{}out.json\\
{}Table 22 & Table typed in the manuscript from the result files of the script & (see the reproducibility appendix of the paper) & Data/\allowbreak{}daisy\_\allowbreak{}ensemble.csv.gz (10 190 runs)\\
{}Figure 19 & Figure drawn by the script from the input data & Manuscripts/\allowbreak{}reproducibility/\allowbreak{}analysis/\allowbreak{}p2\_\allowbreak{}fig6.py;  p2\_\allowbreak{}laws2.py;  p2\_\allowbreak{}laws2.js & Data/\allowbreak{}daisy\_\allowbreak{}ensemble.csv.gz (10 190 runs);  Manuscripts/\allowbreak{}figures/\allowbreak{}p2/\allowbreak{}fig06\_\allowbreak{}compensation.pdf\\
{}Table 23 & Table typed from the result file (19 of its numbers found there) & Manuscripts/\allowbreak{}reproducibility/\allowbreak{}run\_\allowbreak{}all.sh;  Manuscripts/\allowbreak{}reproducibility/\allowbreak{}analysis/\allowbreak{}p2\_\allowbreak{}fig6.py;  Manuscripts/\allowbreak{}reproducibility/\allowbreak{}analysis/\allowbreak{}p2\_\allowbreak{}laws2.py & Data/\allowbreak{}daisy\_\allowbreak{}ensemble.csv.gz (10 190 runs);  Manuscripts/\allowbreak{}results/\allowbreak{}p2\_\allowbreak{}laws2.json\\
{}Table 24 & Table typed from the result file (24 of its numbers found there) & Manuscripts/\allowbreak{}reproducibility/\allowbreak{}run\_\allowbreak{}all.sh;  Manuscripts/\allowbreak{}reproducibility/\allowbreak{}make\_\allowbreak{}e12.py;  Manuscripts/\allowbreak{}reproducibility/\allowbreak{}audit\_\allowbreak{}numbers.py & Data/\allowbreak{}daisy\_\allowbreak{}ensemble.csv.gz (10 190 runs);  Manuscripts/\allowbreak{}results/\allowbreak{}e12.json\\
{}Figure 20 & Figure drawn by the script from the input data & Manuscripts/\allowbreak{}reproducibility/\allowbreak{}analysis/\allowbreak{}p1\_\allowbreak{}fig14\_\allowbreak{}ens.py & Data/\allowbreak{}daisy\_\allowbreak{}ensemble.csv.gz (10 190 runs);  Manuscripts/\allowbreak{}figures/\allowbreak{}p1/\allowbreak{}fig14\_\allowbreak{}ensemble\_\allowbreak{}structure.pdf\\
{}Figure 21 & Figure drawn by the script from the input data & Manuscripts/\allowbreak{}reproducibility/\allowbreak{}analysis/\allowbreak{}p1\_\allowbreak{}fig13\_\allowbreak{}time.py & Data/\allowbreak{}daisy\_\allowbreak{}ensemble.csv.gz (10 190 runs);  Manuscripts/\allowbreak{}figures/\allowbreak{}p1/\allowbreak{}fig13\_\allowbreak{}time\_\allowbreak{}management.pdf\\
{}Figure 22 & Figure drawn by the script from the input data & Manuscripts/\allowbreak{}reproducibility/\allowbreak{}run\_\allowbreak{}all.sh;  Manuscripts/\allowbreak{}reproducibility/\allowbreak{}analysis/\allowbreak{}p1\_\allowbreak{}fig12\_\allowbreak{}drivers.py & Data/\allowbreak{}daisy\_\allowbreak{}ensemble.csv.gz (10 190 runs);  Manuscripts/\allowbreak{}figures/\allowbreak{}p1/\allowbreak{}fig12\_\allowbreak{}drivers.pdf\\
{}Table 25 & Descriptive table (definitions, design, parameters or tests); no computation & - & \\
{}Table 26 & Table typed from the result file (8 of its numbers found there) & Manuscripts/\allowbreak{}reproducibility/\allowbreak{}audit\_\allowbreak{}numbers.py;  Manuscripts/\allowbreak{}reproducibility/\allowbreak{}p3\_\allowbreak{}js/\allowbreak{}e1\_\allowbreak{}heldout.js;  Manuscripts/\allowbreak{}reproducibility/\allowbreak{}analysis/\allowbreak{}p3\_\allowbreak{}figs.py & Data/\allowbreak{}daisy\_\allowbreak{}ensemble.csv.gz (10 190 runs);  Manuscripts/\allowbreak{}results/\allowbreak{}e12\_\allowbreak{}pred.json\\
{}Figure 23 & Figure drawn by the script from the input data & Manuscripts/\allowbreak{}reproducibility/\allowbreak{}fig\_\allowbreak{}vg\_\allowbreak{}check.py & Data/\allowbreak{}daisy\_\allowbreak{}ensemble.csv.gz (10 190 runs);  Manuscripts/\allowbreak{}figures/\allowbreak{}p1/\allowbreak{}fig15\_\allowbreak{}vg\_\allowbreak{}check.pdf\\
{}Table 27 & Table typed from the result file (9 of its numbers found there) & Manuscripts/\allowbreak{}reproducibility/\allowbreak{}p3\_\allowbreak{}js/\allowbreak{}run\_\allowbreak{}drip.js;  Manuscripts/\allowbreak{}reproducibility/\allowbreak{}analysis/\allowbreak{}p3\_\allowbreak{}lit\_\allowbreak{}figs.py;  Manuscripts/\allowbreak{}reproducibility/\allowbreak{}analysis/\allowbreak{}make\_\allowbreak{}supplement\_\allowbreak{}data.py & Data/\allowbreak{}daisy\_\allowbreak{}ensemble.csv.gz (10 190 runs);  Manuscripts/\allowbreak{}results/\allowbreak{}drip\_\allowbreak{}out.json\\
{}Figure 24 & Figure drawn by the script from the input data & Manuscripts/\allowbreak{}reproducibility/\allowbreak{}analysis/\allowbreak{}p1\_\allowbreak{}figs\_\allowbreak{}c.py & Data/\allowbreak{}daisy\_\allowbreak{}ensemble.csv.gz (10 190 runs);  Manuscripts/\allowbreak{}figures/\allowbreak{}p1/\allowbreak{}fig09\_\allowbreak{}recoverable.pdf\\
{}Table 28 & Table typed in the manuscript from the result files of the script & (see the reproducibility appendix of the paper) & Data/\allowbreak{}daisy\_\allowbreak{}ensemble.csv.gz (10 190 runs)\\
{}Table 29 & Table typed from the result file (34 of its numbers found there) & Manuscripts/\allowbreak{}reproducibility/\allowbreak{}run\_\allowbreak{}all.sh;  Manuscripts/\allowbreak{}reproducibility/\allowbreak{}audit\_\allowbreak{}numbers.py;  Manuscripts/\allowbreak{}reproducibility/\allowbreak{}analysis/\allowbreak{}p1\_\allowbreak{}core.py & Data/\allowbreak{}daisy\_\allowbreak{}ensemble.csv.gz (10 190 runs);  Manuscripts/\allowbreak{}results/\allowbreak{}p1\_\allowbreak{}core.json\\
\bottomrule\end{longtable}}


\bibliography{refs}
\end{document}
